\documentclass[letterpaper, 12pt]{book}


\usepackage[utf8]{inputenc}
\usepackage[margin=0.7in]{geometry}
\usepackage[T1]{fontenc}
\usepackage{graphicx}
\usepackage{longtable}
\usepackage{fancyhdr}
\usepackage{wrapfig}
\usepackage{rotating}
\usepackage[normalem]{ulem}
\usepackage{amsmath}
\usepackage{amssymb}
\usepackage{capt-of}
\usepackage{hyperref}
\usepackage{tikz}
\usepackage{blindtext}
\usepackage{subcaption}
\usepackage{titling}
\usepackage{amsmath}
\usepackage{titlesec}
\usepackage{amsmath, amsfonts, amssymb}
\usepackage{multicol}
\usepackage{gensymb}
\usepackage{array}
\usepackage{cancel}
\usepackage{siunitx}
\usepackage{colortbl,hhline}
\usepackage[numbers]{natbib}
\usepackage{indentfirst}
\usepackage{graphicx}
\usepackage{siunitx}
\usepackage[
    backend=bibtex,
    bibencoding=ascii,
    style=chem-angew,   
    citestyle=numeric-comp,
    sorting=none]
{biblatex}

\usetikzlibrary{arrows}
\usepackage{tikz}
\usepackage{pgfplots}
\usepackage[controls]{animate}
\usepackage{color}
\usepackage{epigraph}
\usepackage{empheq}
\usepackage{svg}
\usepackage{mathrsfs}
\usepackage[theorems,skins]{tcolorbox}
\usepackage[svgnames]{xcolor}
\usepackage[explicit]{titlesec}
\usepackage{framed}
\usepackage{centernot}
\usepackage{amsmath}
\usepackage{abstract}
\renewcommand{\abstractname}{}
\renewcommand{\absnamepos}{empty}
\newtheorem{theorem}{Teorema}
\definecolor{shadecolor}{gray}{0.94}
\usepackage{minted}

% Customize section headers
% \titleformat{\section}{\normalfont\sffamily\huge\bfseries\color{blue}}{\thesection}{1em}{}
% \titlespacing{\section}{0em}{.5em}{.5em}
% \titleformat{\subsection}{\large\bfseries\normalfont\sffamily\Large\color{red}}{\thesubsection}{1.5em}{}
% \titlespacing{\subsection}{0em}{.5em}{.5em}

% Customize section and chapter titles
\usepackage{titlesec}
\titleformat{\chapter}[display]{\normalfont\huge\bfseries}{\chaptertitlename\ \thechapter:}{5pt} {\Huge}[\titlerule]

\titlespacing*{\chapter}{0pt}{-20pt}{20pt}


\titleformat{\section}{\normalfont\Large\bfseries}{\thesection}{1em}{} % Section title formatting

\titleformat{\subsection}{\normalfont\large\bfseries}{\thesubsection}{1em}{} % Subsection title formatting
% Add minitoc package
\usepackage{minitoc}
\usepackage[brazil]{babel}

\usepackage{pifont}
\usepackage{amssymb}
\usepackage[T1]{fontenc} %Not needed by LuaLaTeX or XeLaTeX
\usepackage[brazilian]{babel}
\usepackage[utf8]{inputenc}
%%% Local Variables:
%%% mode: latex
%%% coding: utf-8
%%% TeX-master: "main.tex"
%%% End:




\addbibresource{refs.bib}
% \author{Jefter Santiago}
\date{}


\usepackage[utf8]{inputenc}
\usepackage[margin=0.7in]{geometry}
\usepackage[T1]{fontenc}
\usepackage{graphicx}
\usepackage{longtable}
\usepackage{fancyhdr}
\usepackage{wrapfig}
\usepackage{rotating}
\usepackage[normalem]{ulem}
\usepackage{amsmath}
\usepackage{amssymb}
\usepackage{capt-of}
\usepackage{hyperref}
\usepackage{tikz}
\usepackage{blindtext}
\usepackage{subcaption}
\usepackage{titling}
\usepackage{amsmath}
\usepackage{titlesec}
\usepackage{amsmath, amsfonts, amssymb}
\usepackage{multicol}
\usepackage{gensymb}
\usepackage{array}
\usepackage{cancel}
\usepackage{siunitx}
\usepackage{colortbl,hhline}
\usepackage[numbers]{natbib}
\usepackage{indentfirst}
\usepackage{graphicx}
\usepackage{siunitx}
\usepackage[
    backend=bibtex,
    bibencoding=ascii,
    style=chem-angew,   
    citestyle=numeric-comp,
    sorting=none]
{biblatex}

\usetikzlibrary{arrows}
\usepackage{tikz}
\usepackage{pgfplots}
\usepackage[controls]{animate}
\usepackage{color}
\usepackage{epigraph}
\usepackage{empheq}
\usepackage{svg}
\usepackage{mathrsfs}
\usepackage[theorems,skins]{tcolorbox}
\usepackage[svgnames]{xcolor}
\usepackage[explicit]{titlesec}
\usepackage{framed}
\usepackage{centernot}
\usepackage{amsmath}
\usepackage{abstract}
\renewcommand{\abstractname}{}
\renewcommand{\absnamepos}{empty}
\newtheorem{theorem}{Teorema}
\definecolor{shadecolor}{gray}{0.94}
\usepackage{minted}

% Customize section headers
% \titleformat{\section}{\normalfont\sffamily\huge\bfseries\color{blue}}{\thesection}{1em}{}
% \titlespacing{\section}{0em}{.5em}{.5em}
% \titleformat{\subsection}{\large\bfseries\normalfont\sffamily\Large\color{red}}{\thesubsection}{1.5em}{}
% \titlespacing{\subsection}{0em}{.5em}{.5em}

% Customize section and chapter titles
\usepackage{titlesec}
\titleformat{\chapter}[display]{\normalfont\huge\bfseries}{\chaptertitlename\ \thechapter:}{5pt} {\Huge}[\titlerule]

\titlespacing*{\chapter}{0pt}{-20pt}{20pt}


\titleformat{\section}{\normalfont\Large\bfseries}{\thesection}{1em}{} % Section title formatting

\titleformat{\subsection}{\normalfont\large\bfseries}{\thesubsection}{1em}{} % Subsection title formatting
% Add minitoc package
\usepackage{minitoc}
\usepackage[brazil]{babel}

\usepackage{pifont}
\usepackage{amssymb}
\usepackage[T1]{fontenc} %Not needed by LuaLaTeX or XeLaTeX
\usepackage[brazilian]{babel}
\usepackage[utf8]{inputenc}
%%% Local Variables:
%%% mode: latex
%%% coding: utf-8
%%% TeX-master: "main.tex"
%%% End:



% Document config
\hypersetup{
  pdfauthor={Jefter Santiago},
  pdftitle={Anotações de IntroFismat},
%  pdfkeywords={},
  pdfsubject={Introdução à Física Matemática},
  pdfcreator={Jefter Santiago}, 
  pdflang={Portuguese},
  colorlinks=true,    % Color links instead of boxes
  linkcolor=blue,     % Color of internal links
  citecolor=green,    % Color of citation links
  urlcolor=blue,      % Color of URLs
}

% remove paginas vazias do inicio do documento
\let\cleardoublepage=\clearpage

\newgeometry{a4paper, left=1.5cm, right=1.5cm, top=2.5cm, bottom=2.5cm}
\begin{titlepage}
  \vspace*{\fill}
  \begin{center}
    \Huge{Introdução à Física Matemática (7600016)}
  \end{center}
  \begin{center}
    \Large{Anotações, exercícios etc}
  \end{center}
  \begin{center}
    \large{Jefter Santiago}
  \end{center}
  \vspace*{\fill}
\end{titlepage}
%%% Local Variables:
%%% mode: latex
%%% coding: utf-8
%%% TeX-master: "main.tex"
%%% End:

\begin{document}
\maketitle
\setcounter{tocdepth}{1} \tableofcontents
\dominitoc\setcounter{tocdepth}{1}

\chapter*{Prefácio - Descrição dos conteúdos das anotações}
Essas são minhas anotações feitas no decorrer da disciplina de Introdução à Física Matemática no
IFSC-USP em 2023.
A ementa da disciplina pode ser encontrada no site oficial da USP, o
\href{https://uspdigital.usp.br/jupiterweb/obterDisciplina?sgldis=7600016&codcur=76041&codhab=0}{jupiter}.
De qualquer forma coloquei os conteúdos ``previstos'' pela ementa abaixo.
\begin{itemize}
\item[\rlap{\raisebox{0.3ex}{\hspace{0.4ex}\tiny \ding{52}}}$\square$] Números complexos e suas propriedades: interpretações geométricas (cartesiana, polar e vetorial), fórmulas de Euler e de De Moivre. 
\item[\rlap{\raisebox{0.3ex}{\hspace{0.4ex}\tiny \ding{52}}}$\square$] Funções complexas de um parâmetro real: diferenciação e integração. 
\item[\rlap{\raisebox{0.3ex}{\hspace{0.4ex}\tiny \ding{52}}}$\square$] Equações diferenciais ordinárias de primeira ordem: equações lineares, equações separáveis e equações exatas, fator integrante. 
\item[\rlap{\raisebox{0.3ex}{\hspace{0.4ex}\tiny \ding{52}}}$\square$] Equações diferenciais ordinárias de segunda ordem: soluções linearmente independentes, Wronskiano, método da variação das constantes. 
\item[\rlap{\raisebox{0.3ex}{\hspace{0.4ex}\tiny \ding{52}}}$\square$] Sistemas de equações lineares de primeira ordem: autovalores e autovetores de matrizes, matriz fundamental de sistemas com coeficientes constantes. 
\item[\rlap{\raisebox{0.3ex}{\hspace{0.4ex}\tiny \ding{52}}}$\square$] Transformada de Laplace: propriedades básicas, funções descontínuas e impulsivas (delta de Dirac), relações com a função gama e equações integrais, conceito de convolução. 
\item[\rlap{\raisebox{0.3ex}{\hspace{0.4ex}\tiny \ding{52}}}$\square$] Soluções em séries de equações de segunda ordem: pontos ordinários e singulares regulares (método de Frobenius), ênfase nos polinômios de Legendre.
\end{itemize}

As anotações são baseadas nos que foi escrito em sala pelo professor Vanzella e nos livros/anotações
de aula, que usei para estudar. Essas referências estão listadas no fim do documento.

% \chapter{Números Complexos}
%\section{Álgebra dos números complexos}
\blindtext
\blindtext
\blindtext

\section{Representação em forma polar - Fórmula de Euler}
\blindtext
\blindtext
\blindtext
\blindtext

\section{Funções complexas elementares}
\section{Funções complexas inversas}

%%% Local Variables:
%%% mode: latex
%%% coding: utf-8
%%% TeX-master: "main.tex"
%%% End:




%\chapter{Equações diferenciais}
%Exemplos de equações diferenciais:

\paragraph{(1) Sistema com particula em queda livre}

\paragraph{(2) Sistema com particula sofrendo força de resistência}

\paragraph{(3) Oscilador harmônico}


\subsection{Classificação de EDOs}


\subsection{Campo de direções}


%%% Local Variables:
%%% mode: latex
%%% coding: utf-8
%%% TeX-master: "main.tex"
%%% End:


\chapter{EDOs de primeira ordem}
% \begin{wrapfigure}{r}{0.3\textwidth}
%   \centering
%   \includegraphics[width=0.3\textwidth]{./images/xkcd/dPain_over_dt.png}
% \end{wrapfigure}

\section{EDOs homogêneas, inomogêneas - fator integrante}
Para equações homogêneas e inomogêneas lineares com variável
independente \( t \) e dependente \( y \), queremos saber a
cara de \( y(t) \) dada na equação \( a(t) y'(t) + b(t) y(t) = c(t) \) 

onde \( a, b  \) e \( c \) são funções dadas de \( t \) .


\paragraph{$\blacksquare$ Exemplo}: \( t^2 y' + 5t y = t^3 - 1 \)
Se \( c(t) = 0 \) então a EDO é homogênea e caso contrário inomogênea.

\begin{itemize}
  \item[$\blacksquare$] Caso homogêneo: 
  \[ a(t) y'(t) + b(t) y(t) = 0 \]
  \[ y'(t) + p(t) y(t) = 0, p(t) = \frac{b(t)}{a(t)} \]
  \[ \Leftrightarrow \frac{y'(t)}{y(t)} = - p(t) \Leftrightarrow \frac{1}{y} \frac{dy}{dt} = -p(t) \]
  \[ \Leftrightarrow \int \frac{1}{y} dy = - \int p(t) dt \]
  \[ \Leftrightarrow  \int \frac{dy}{y} = - \int p(t) dt  \]
  \[ \Leftrightarrow  \ln |y| = - \int p(t) dt \]
  \[ \Leftrightarrow  |y| = e^{- \int p(t) dt + \mathbb{C}} = e^{-\int p(t) dt} e^{\mathbb{C}} \]
  \[ \Leftrightarrow  |y| = |A|e^{- \int p(t) dt}, |A| = e^{\mathbb{C}} \]
  \[ \Leftrightarrow  y(t) = Ae^{- \int p(t) dt}\]

  O valor de $A$ vai depender das condições iniciais do problema.
  
  \item[$\blacksquare$] Caso inomogêneo: \textbf{fator integrante}

  \[ a(t) y'(t) + b(t) y(t) = c(t) \]
  \[ y'(t) + p(t) y(t) = r(t), p(t) = \frac{b(t)}{a(t)}, r(t) = \frac{c(t)}{a(t)} \]

  A estratégia para resolver com fator integrante é chegar à uma equação
  que permita-nos escrever o lado esquerdo da equação usando a regra do produto
  para derivadas, ou seja,
  \[ \frac{d  }{d x} \left( f g \right) = g \frac{d}{dx}f  + f \frac{d}{dx}g \]

  Para tal, precisamos de uma combinação linear do tipo
  \( \mu (t) y'(t) + \mu (t) p(t) y(t) = \mu (t) r(t) \), onde \( \mu(t) \) é o fator
  integrante.
  A relação
  \[ \frac{d}{dt} \left(\mu(t) y(t) \right) = \mu(t) y'(t) + \mu(t) p(t) y(t) \] 

  é válida so e se só \( \mu(t) \) obedece a relação

  \begin{equation}
    \frac{d}{dt} \mu(t) = p(t) \mu(t) 
    \label{eq:fator_integrante}
  \end{equation}
\end{itemize}


\paragraph{$\blacksquare$ Exemplo}: \( t^2 y'(t) + 5ty(t) = t^3 - 1 \) 

\[ \Leftrightarrow y' + \frac{5}{t}y = t - \frac{1}{t^2}  \]
\[ \Leftrightarrow \frac{d}{dt}\mu = \mu y' + \frac{5}{t}\mu y = \left( t - \frac{1}{t^2}\right)\mu \]
\[ \Leftrightarrow  \frac{d\mu}{dt}= \frac{5\mu}{t} \]
Essa relação é uma EDO homogênea simples, então temos
\[ \Leftrightarrow \frac{\mu'}{\mu} = \frac{5}{t} \]
\[ \Leftrightarrow \ln|\mu| = 5 \ln |t| \Rightarrow \mu = t^5\]

voltando à equação anterior temos

\[ \frac{d}{dt} \left( t^5 y \right) = t^6 - t^3 \]
\[ t^5 y = \frac{t^7}{7} - \frac{t^4}{4} + \mathbb{C} \]

\[ \Rightarrow y(t) = \frac{t^2}{7}- \frac{1}{4t} + \frac{\mathbb{C}}{t^5} \]


\section{Método de variação de paramêtros}
A solução de uma EDO linear inomogênea de primeira ordem é a soma da solução geral da homogênea
associada com uma solução particular \textbf{qualquer} da EDO inomogênea original.
Podemos portanto `` construir '' a solução de uma EDO inomogênea.


\[ y' + p(t) y = r(t) \rightarrow y(t) = y_{H} (t) + y_{P} (t) \]
onde \( y_{H} \) é a solução da EDO homogênea associada, \( y_P \) é a paramétrica obtida a partir
da homogênea e \( y \) é a solução geral da EDO dada.

\paragraph{$\blacksquare$ Exemplo} \( t y' - t^2 = y, t >0 \) 

\[ t y' - y = t^2 \Rightarrow y' - \frac{1}{t}y = t  \]
Fator integrante \( \mu y' - \frac{\mu}{t}y = \mu t \Rightarrow \frac{d}{dt} \left( \mu y \right) = \mu t \) 
\[ \mu ' = - \frac{\mu }{t} \Leftrightarrow \frac{\mu'}{\mu} = - \frac{1}{t} \Rightarrow \frac{d\mu}{\mu}= - \frac{dt}{t} \Leftrightarrow \ln|\mu| = - \ln|t|\]
\[ \Rightarrow \mu = \frac{A}{t} \]

\[\Rightarrow \frac{d}{dt} \left( \frac{y}{t} \right)  = 1 \Leftrightarrow \frac{1}{t} y = 1 \Leftrightarrow \frac{1}{t}y = t +
  \mathbb{C} \]
\[ \Leftrightarrow y(t) = t^2 + \mathbb{C} t \]

Para \( y(1) = 0 \Rightarrow 0 = 1^2 + \mathbb{C} \cdot 1 \Rightarrow \mathbb{C} = -1 \Leftrightarrow y(t) = t^2 - t \) 

Checando a propriedade da variação de paramêtros a partir desse resultado.
Chutando \( y_P (t) = t^2 \) a solução geral da homogênea: \( y_{H}' - \frac{1}{t} y_H = 0 \) 
\[ \frac{y'_{H}}{y_{H}} = \frac{1}{t} \Leftrightarrow \frac{d y_{H}}{y_{H}} = \frac{dt}{t} \Leftrightarrow \ln|y_H| = \ln t + \mathbb{C}  \]
\[  |y_H(t)| = e^{\mathbb{C}} t \Leftrightarrow y_H(t) = \mathbb{C} t \]
\[ \Leftrightarrow y(t) = \mathbb{C} t + t^2  \]

\section{EDOs separáveis}
Uma EDO de primeira ordem é dita ser separável se pode ser colocada na forma

\begin{equation}
  f(y) y' = g(t) \leftrightarrow M(y)y' + N(t) = 0 \Leftrightarrow M(y) dy + N(t) dt = 0
  \label{eq:edo_separavel}
\end{equation}

\section{Modelagem de sistemas usando EDOs}
Tratamos as populações a partir de densidades para evitar trabalhar com sistemas discretos que podem
levar a descontinuidades. Na realidade quando implementamos alguma simulação podemos tentar
considerar as discretizações do sistema. Na teoria estamos tratando dos modelos mais simples
possíveis.
\subsection{EDOs autonômas e crescimento populacional}
Sejam
\begin{itemize}
 \item[$\bullet$] y: densidade de ``entes'': Substâncias, presas $\times$ predadores, moleculas...
 \item[$\bullet$] t: tempo
\end{itemize}

Supondo \( \frac{dy}{dt} \propto y \Leftrightarrow \frac{dy}{dt}=  r y, r > 0 \) , onde \( r \) é a taxa de
crescimento. Esse é o modelo bem simplificado de um sistema fechado.

Temos
\[ \frac{dy}{y}= r dt \Leftrightarrow \ln | y| = rt + \mathbb{C} \Leftrightarrow y(t) = A e^{rt}\]
supondo \( y(t_0) = y_0 \Rightarrow y_0 = A e^{rt_0} \Rightarrow A = y_0 e^{-rt_0} \), logo
\[ y(t) = y_0 e^{r(t-t_0)} \]
note que para qualquer \( t \) a solução é apenas transladada.

\subsection{Equação logística}
Leva em consideração um limite superior para a solução. Isso é mais ``realista'' já que tem um fator
limitante como população cujo crescimento é proporcional.

\[ \frac{dy}{dt}= r(y) y = r_0 \left( 1 - \frac{y}{k} \right)y, r(y) \text{ é decrescente. }\]

[IMAGEM] 

Ponto \( k \) (valor crítico): a partir desse ponto a taxa de ``mortalidade'' é menor que a
de ``natalidade'', ou seja, a população decresce.

\[ f(y) = \frac{dy}{dt} = r_0 \left( 1 - \frac{y}{k} \right)y \]

[IMAGEM] [IMAGEM]

\subsection{Resolvendo a equação logística}
\[ \frac{dy}{dt} = r_0 \left( 1 - \frac{y}{k} \right)y  \Leftrightarrow y(t) = 0, y(t) = k \]
são soluções, logo \( \frac{dy}{y \left( 1 - \frac{y}{k} \right)}  = r_0 dt \)
integrar essa relação nos leva à uma integral por fração parcial. A forma de resolver
isso está no apendice [citar aqui]

resolvendo a integral temos

\[ \ln \bigg| \frac{y}{y- k} \bigg | = r_0 t + \mathbb{C} \]
\[ \frac{y}{y-k} = A e^{r_0 t} \Rightarrow y = A e^{r_0 t} (y - k) \]
\[ y \left( A e^{r_0 t} \right) = - k e^{r_0 t} A \]
\[ \Rightarrow y(t) = \frac{K}{-A^{-1} e^{-r_0 t} + 1} \]

\section{EDOs exatas}

Partindo das EDOS separáveis, temos \( f(y)dy + g(t) dt = 0 \),
ou seja \( F(y) + G(t) = \text{cte}, \frac{dF}{dy} = f(y), \frac{dG}{dt} = g(t) \) 

\paragraph{Generalizando}: supondo que existem um classe de EDOS cujas soluções podem ser dadsa por
 \( \psi (t, y(t)) = \text{cte} \), onde \( \psi(t, y(t)) = F(t, y(t)) + G(t, y(t)) \). A função
 \( \psi \) é um lugar geométrico em \( x, y \), uma curva de nível no plano de forma geral
 \( x = t \)  e \( y = y(t) \), para \( \psi(x, y) \).
 \[ \left( \frac{\partial }{\partial x} \psi \right) \frac{dx}{dt} + \left( \frac{\partial }{\partial y} \psi \right)\frac{d y }{d t} = 0 \]

 Com \( \psi_{, x} = \frac{\partial }{\partial x} \psi, \psi_{, y} = \frac{\partial }{\partial y} \psi\), substituindo na equação acima temos
 \[ \psi_{,x} + \psi_{,y}y' = 0 \]
 Isso é uma EDO separáel e pode ser escrita como
 \[ M(x, y) + N(x, y) y'(x) = 0 \]
 onde \( M(x, y) = \frac{\partial }{\partial x} \psi \)  e \( N(x, y) = \frac{\partial }{\partial y} \psi \), portanto
 \[ \psi(x, y(x)) = \text{cte} \]
 
 Qualquer EDO que pode ser escrita dessa forma é chamada de EDO exata.



 \paragraph{$\blacksquare$} Explicação (rasa) sobre essa manipulação:

 \[ \nabla \psi = \left( \frac{\partial }{\partial x} \psi, \frac{\partial }{\partial y} \psi \right) \]
 identificação de \( \frac{\partial }{\partial x} \psi \)  e \( \frac{\partial }{\partial y} \psi \), vale que
 \( \nabla \times \psi = 0 \Leftrightarrow \nabla \times \psi = \frac{\partial }{\partial y} \psi - \frac{\partial }{\partial x} \psi = 0  \Leftrightarrow \frac{\partial M}{\partial y} = \frac{\partial N}{\partial x}\).
 Essa é a condição necessária (e nesse caso suficiente) para encontrar a função \( \psi \).
 



    \subsection{Caso base}
    \paragraph{$\blacksquare$ Exemplo}: $(\sin x + x^2 e^y -1) y' + (y \cos x + 2x e^y ) = 0$ 
    Verificar validade de \( M(x, y) + N(x, y) y' = 0 \), com \( M_{,x} \overset{?}{=} N_{,y} \) 

    \[ M(x, y) = y \cos x + 2x e^y \mapsto \frac{\partial }{\partial y} M = \cos x + 2x e^y \]
    \[ N(x, y) = \sin x + x^2 e^y - 1 \mapsto \frac{\partial }{\partial x} N = \cos x + ex e^y \]

    \( M_{,x} = N_{,y} \), a equação é exata.


    Encontrando a função \( \psi(x, y) \) 

    \[ \frac{\partial \psi}{\partial x}  = y \cos x + 2x e^y \Rightarrow \psi(x, y) = y \sin x + x^2 e^y + h(y) \]
    fazendo \( \frac{\partial }{\partial y} \psi(x, y) = \frac{\partial }{\partial y}  \left( y \sin x + x^2 e^y + h(y) \right) \) 
    \[ \frac{\partial }{\partial y} \psi(x, y) = \sin x + x^2 e^y + h'(y) = \sin x + x^2 e^y - 1 \]
    \[ \Leftrightarrow h'(y) = -1 \Leftrightarrow h(y) = -y + \mathbb{C} \]
    \[ \Longrightarrow  \psi(x, y) = y \sin x + x^2 e^y - y + \mathbb{C} \]

    Portanto, a solução é dada pelo \( y(x) \) que satistfaz \( y(x) \sin x + x^2 e^{y(x)} - y(x) = \text{cte} \).



    \subsection{Caso que precisa de fator integrante}
    
 
    \paragraph{$\blacksquare$ Exemplo}: equação não exata e fator integrante para equações não exatas
    \[(2 x + 4y) + (2x - 2y) y' = 0\]


    Verificar a validade \( M(x, y) + N(x, y) y' = 0 \):
    \[ \begin{cases}
         M = 2x + 4y \\
         N = 2x - 2y
       \end{cases} \Rightarrow \begin{cases}
                       M_{, y} = 4 \\
                       N_{, x} = 2
                     \end{cases} \Rightarrow \text{Não satisfazem a condição} \]

     Para resolver usamos fator integrante: supondo que \( M(x, y) + N(x, y)y' = 0 \) é não
     exata. Encontrar \( \mu = \mu(x, y) \), tal que \( \mu M + \mu N y' = 0 \) seja exata.

     \[ \widetilde{M} = \mu M , \widetilde{N} = \mu N \]
     então \( \widetilde{M} + \widetilde{N} y' = 0 \), logo
     \[ (\mu M)_{, y} = (\mu N)_{,x} \]
   



     \paragraph{$\blacksquare$ Exemplo}: $ \left( 2 \sin(y) - x \right) y' = \tan(y)$
     \[ M(x, y) = - \tan(y) \Rightarrow \frac{\partial }{\partial y} M(x, y) =  -  \frac{\partial }{\partial y} \left(
         \frac{\sin(y)}{\cos(y)} \right) = - \frac{\cos (y)}{\cos(y)} - \frac{\sin^2 (y)}{\cos^2
         (y)}= - 1 - \tan^2 (y) = - \sec^2 (y) \]
    

     \[ N(x, y) = 2 \sin(y) - x \Rightarrow \frac{\partial }{\partial x}  N(x, y) = -1  \]
     
     \[ N_x  \neq  M_y \]
     
     Logo, pela forma que foi dada não é uma equação exata, precisamos encontrar algum fator
     integrante para resolve-la.
     
     Nessa situação chutamos primeiro fatores integrantes do tipo \( \mu(x, y) \) e outras variações
     que envolvam as duas variáveis ou uma delas, como \( \mu(x) \), \( \mu(y) \), \( \mu(x/y) \), \(
     \mu(xy) \), \( [\reflectbox{$\dots$}] \). A partir do fator integrante substituimos e checamos se
     a equivalência se permanece. Vimos isso na sequência.


     Supondo que para essa EDO podemos usar o fator integrante \( \mu = \mu(x, y) \), então
     \( \widetilde{M} + \widetilde{N} y'   \mu M + \mu N y' = 0 \) 
     \[ \mu(x, y) \left[ 2 \sin (x) - x \right]y' + \mu(x, y) (-\tan(y)) = 0 \]
     
     Checando a validade de \( \frac{\partial }{\partial y} \widetilde{M}  = \frac{\partial }{\partial x} \widetilde{N}\):
     
     
     \[ - \frac{\partial \mu}{\partial y} \tan(y) - \mu \sec^2 (y) = \frac{\partial \mu}{\partial x} \left( 2 \sin (y) - x \right) - \mu (-1) \]
     \[ \Leftrightarrow \left( x - 2 \sin(y) \right) \frac{\partial \mu}{\partial x} - \tan(y) \frac{\partial \mu}{\partial x} = (\sec^2 (y) - 1) \mu  \]
     
     \paragraph{$\rightarrow $ Ansatz}: \( \mu = \mu(x)\Rightarrow \frac{\partial \mu}{\partial x} = \mu' \)  e \( \frac{\partial \mu}{\partial y} = 0 \) 
     
     \[ \Leftrightarrow (x - 2\sin(y)) \mu'(x) = \left( \tan(y) \right)^2 \mu(x) \]
     \[ \Leftrightarrow \frac{\mu'(x)}{\mu(x)} = \frac{ \left( \tan(y) \right)^2 }{x - 2 \sin(y)} \]

     Essa última linha é uma contradição, já que o lado direito não depende explicitamente apenas de
     \( x \), ou seja, esse chute não é válido.
     

     \paragraph{$\rightarrow $ Ansatz}: \( \mu = \mu(y) \Rightarrow \frac{\partial \mu}{\partial x}  = 0, \frac{\partial \mu}{\partial y} = \mu'(y)\) 

     \[ - \tan(y) \mu' = \left( \tan(y) \right)^2 \mu \Rightarrow \frac{\mu'(y)}{\mu(y)}= - \tan(y) \Rightarrow \mu(y) = \cos(y)\]



     \paragraph{$\blacksquare$ Exemplo}: $(2 \sin(y) \cos(y) - x \cos (y))y' - \sin(y) = 0$

     \[ M_y = - \cos (y) , N_{x} = - \cos(y) \]

     \[ M = \psi_{x} = - \sin (y) \Rightarrow \psi (x, y) = - x \sin (y) + h(y)\]
     \[ N = \psi_{y} = \sin(y) - x \cos(y) = - x \cos(y) + h'(y) \Rightarrow h'(y) = \sin(2y) \Rightarrow h(y) = -
       \frac{1}{2} \cos(2y) \]
     \[  \psi(x, y) = - x \sin(y) - \frac{1}{2} \cos(2y) + \mathbb{C} \]

     Logo, a solução implicita da EDO é

     \[ -x \sin(y(x)) - \frac{1}{2} \cos(2 y(x)) = \mathbb{C}  \]

     
%%% Local Variables:
%%% mode: latex
%%% coding: utf-8
%%% TeX-master: "main.tex"
%%% End:



\chapter{EDOs de segunda ordem}

% \begin{wrapfigure}{r}{0.5\textwidth}
%   \centering
%   \includegraphics[width=0.3\textwidth]{./images/xkcd/differentiation_and_integration.png}
% \end{wrapfigure}

\section{Lineares homogêneas}

Os métodos para resolver as EDOs de segunda ordem lineares são os métodos
do Wronskiano (usando o teorema de Abel), redução de ordem, variação de parametros (para
as EDOs inomogêneas) e em último caso, método de séries. Nessa etapa vimos todos métodos exceto esse
último, que será visto na sequência.

As EDOs de segunda ordem tem a seguinte cara:
\[ a(t) y'' + b(t) y' + c(t) y = 0  \]

vamos considerar inicialmente só o caso homogêneo por simplicidade e sempre
adotar o uso da equação na forma ``canônica'', isto é, com a derivada de maior ordem multiplicada
por \( 1 \):

\[ y''(t) + p(t) y'(t) + q(t) y(t) = 0, p(t) = \frac{b(t)}{a(t)}, q(t) = \frac{c(t)}{a(t)} \]


As EDOs de segunda ordem são compostas de 2 soluções linearmente independentes, isso significa que
podemos construir uma solução geral a partir de outra e em geral é isso que vamos fazer, supor uma
possível solução, com base por exemplo na ``cara'' da EDO.Isto vai ser visto melhor na prática em
exemplos desenvolvidos.

Segue propriedades gerais das EDOs de segunda ordem lineares:


\begin{theorem}[Linearidade ou principio da superposição]
  Se \( y_1 (t)  \)  e \( y_2 (t) \) são soluções de uma EDO linear homogênea, então
  \( \alpha y_1 (t) + \beta y_2 (t) \), \( \forall \alpha , \beta  \in \mathbb{R} \), também é solução.
\end{theorem}

\begin{theorem}[Solução geral]  
  Se \( y_1(t) \) e \( y_2 (t) \) são soluções de uma EDO linear homogênea de 2ª ordem que não
  são proporcionais uma à outra, então \( y(t) = \alpha y_1 (t) + \beta y_2 (t) \), com \( \alpha  \) e
  \( \beta  \) constantes quaisquer é a solução geral da EDO dada.
\end{theorem}


\begin{theorem}[Wronskiano]\label{teo:definicao_wronskiano}
  Dadas duas funções quaisquer \( y_1(t) \) e \( y_2(t) \) define-se o chamado
  Wronskiano como sendo a função
  \begin{equation}
    W[y_1, y_2] (t) := y_1 (t) y_2 '(t) - y_2 (t) y_1 ' (t)
    \label{eq:wronskiano_formula}
  \end{equation}
\end{theorem}

\paragraph{Expandindo um pouco a ideia do Wronskiano}

Nota-se que, \( W = 0 \) só e só se \( y_1 (t) = f(t) \) e \( y_2 (t) = c f(t)\) .

Derivando o Wronskiano:
\[ \frac{d W }{d t} = \frac{d  }{d t} \left[ y_1 y_2 ' - y_2 y_1 ' \right] = y_1'y_2 ' + y_1 y_2''
  - y_2' y_1' - y_2 y_1'' \]
\[ \frac{d W }{d t} = y_1 y_2'' - y_2 y_1'' \]

\[ \frac{d W }{d t} = y_1 \left\bigg( -p(t) y_2 ' - q(t) y_2  \right) - y_2 \left\bigg( -p(t) y_1' - q(t) y_1 \right)\]

\[ \frac{d W }{d t} = -p(t) \left\bigg( y_1 y_2 ' - y_2 y_1' \right) = -p(t) W\]

\begin{theorem}[Teorema de Abel]\label{eq:teorema_abel}
  Se \( y_1 \) e \( y_2 \) são soluções da EDO \( y'' + p(t) y' + q(t) y = 0 \) seu Wronskiano
  satisfaz \( W' + p(t) W = 0 \) vale
  \begin{equation}
    W[y_1, y_2] (t) = W_0 \exp\left({- \int p(t) dt}\right)
    \label{eq:teorema_abel}
  \end{equation}
\end{theorem}

\paragraph{$\blacksquare$ Exemplo}: \( x^2 y'' - x(x + 2)y' + (x+2) y = 0, x > 0, y(x) = \)?
Sabemos que a solução geral deve ser da forma \( y(x) = \alpha y_1 (x) + \beta y_2 (x) \) 
\[ y'' - \frac{x(x+2)}{x^2}y' + \frac{(x + 2)}{x^2}y = 0 \]
\[ y'' - \left( 1 + \frac{2}{x} \right)y' + \frac{1}{x} \left( 1 - \frac{2}{x} \right) y = 0 \]

\[ W = y_1 y_2' - y_2 y_1' = W_0 \exp{\int \left( 1 + \frac{2}{x} \right) dx} = W_0 \exp{x + 2 \ln(x)} \]
\[ y_1 y_2' - y_2 y_1' = W_0 x^2 e^x \]

Chutando uma solução para a EDO:
Olhando para a essa EDO podemos chutar uma solução do tipo \( y(x) = x^{\lambda} \), já que derivando duas
vezes vamos fazer o coeficiente \( x^2 \) sumir da EDO. Primeiro fazemos essa substituição na EDO
original e tentar calcular o \( \lambda  \):

\[ \begin{cases*}
     y' = \lambda x^{\lambda - 1} \\
     y'' = \lambda (\lambda -1) x^{\lambda - 2} 
   \end{cases*} \Longrightarrow 
   \begin{align*}
     \lambda (\lambda - 1) x^{\lambda } - \lambda (\lambda + 2) x^{\lambda } + (x+2)x^{\lambda } & = 0  \\
     x^{\lambda} \left[ \lambda (\lambda -1) - 2 \lambda  + 2 \right] + (1 - \lambda ) x^{\lambda + 1} &= 0
   \end{align*}
 \]

 Se \( \lambda  = 1 \) a relação é válida, logo \( y_1(x) = x \) pode ser uma solução.


 A partir desse resultado usamos o (\ref{eq:teorema_abel}) para encontrar a segunda solução. Caimos
 em uma EDO de primeira ordem e resolvendo ela temos \( y_2(x) \), e portanto a solução completa.

 \[ y_1 y_2' - y_2 y_1' = W_0 x^2 e^x \]
 \[ x y_2'- y_2 = W_0 x^2 e^x \Longrightarrow y_2' - \frac{1}{x} W_0 x e^x, W_0 \neq 0\]

 \[ \mu = \frac{1}{x} \Longrightarrow  \frac{d  }{d x} \left( \frac{y_2}{x} \right) = e^x \Longrightarrow  \frac{y_2}{x} = e^x + \mathbb{C} \]
 \[ \Longrightarrow y_2(x) = x e^x \]

 Logo, a solução geral da EDO dada é \( y(x) = \alpha y_1(x) + \beta y_2(x) \) 
 \[ y(x) = \alpha x + \beta x e^x \]

 A partir das condições iniciais podemos definir os coeficientes \( \alpha \) e \( \beta  \).

\subsection{Método de redução de ordem}


Partindo de uma EDO de segunda ordem conseguimos encontrar uma solução a partir de outra
substituindo por uma EDO de primeira ordem.

\[ a y'' + b y' + cy = 0 \]

Supondo que conhecemos uma das soluções \( y_1 \) que obtemos a partir de algum Anzats. Podemos
então obter a \( y_2 \) por redução de ordem fazendo \( y_2 = v(t)y_1(t) \) e substituindo na EDO
original
\[ \begin{cases}
     y_2 ' = v' y_1 + v y_1' \\
     y_2 '' = v'' y_1 + 2v'y_1' + v y_1''
   \end{cases}
   \Longrightarrow
   \begin{align*}
     a \left\bigg( v''y_1+2v'y_1'+vy_1'' \right) + b \left\bigg( v'y_1+vy_1' \right) + c v y_1 & =
                                                                                                 0\\
     v \underbrace{\left[ ay_1'' + by_1' + cy_1 \right]}_{= 0} + av''y_1+2v'y_1'+bv'y_1 & = 0 \\
   \end{align*}
 \]

 \[
 av''y_1 + 2v'y_1' + bv' y_1 = 0 \Longrightarrow  v'' + v' \left\bigg[ \frac{2y_1' + b y_1}{ay_1} \right] = 0
 \]

 a partir desse ponto substituimos o \( y_1, a \) e \( b \), que são conhecidos e resolvemos a EDO
 em \( v'' \) para depois substituir em \( y_2 \). Não continuei essa demonstração porque me parece
 ser mais útil tratar um exemplo real ao invés.


 \paragraph{$\blacksquare$ Exemplo}: \( 4 y'' + 4y' + y = 0 \)
 
 Supondo \( y_1 = e^{r t} \Rightarrow y_1' = ry_1 \Rightarrow y_1'' = r^2 y_1\) e disso segue que \( y \left( 4r^2 + 4r
   + 1 \right) = 0 \Rightarrow (2r + 1)^2 = 0 \Rightarrow r = -\frac{1}{2} \) é uma solução. Aplicamos a redução de
 ordem para encontrar a segunda solução.

 \[ y_2 = v y_1 \Rightarrow \begin{cases}
                    y_2' = v' y_1 + vy_1' \\
                    y_2'' = v''y_1 + 2v'y_1' + vy_1''
                  \end{cases}  \]

       
  \[ \left( v'' y_1 + 2 v' y_1' + vy_1'' \right)  + p  () + q \left( vy_1 \right)  =  0\]                
  \[ \Rightarrow v'' y_1 + v' \left( 2y_1' + p y_1 \right) = 0  \]

  Resolvendo essa EDO para encontrar \( y_2 \) :
  \[ y_1' = - \frac{1}{2} y1 , y_1'' = \frac{1}{4} y_1 \]
  \[ y_2 = vy_1 \Rightarrow \begin{cases}
                    y_2' = v' y_1 - \frac{v}{2}y_1'\\
                    y_2'' = v''y_1 - \underbrace{\frac{v'y_1}{2}- \frac{vy_1'}{2}}_{= - v y_1} + \frac{vy_1}{4}
                  \end{cases} \]

  Substituindo na equação original:
                
\[ 4 \left\right( v'' y_1 - v y_1 + \frac{v}{4}y_1 \right) + 4 \left\bigg( v' y_1 - \frac{v}{2}y_1
\right) + v y_1 = 0 \] 
\[ \Rightarrow 4 v'' y_1 = 0 \Rightarrow v'' = 0 \Rightarrow v' = a \Rightarrow v = a t + b \]

\[ y_2 = v y_1 = at e^{-t/2} + b e^{-t/2} \]
                
\section{Lineares e inomogêneas}
\subsection{Solução com coeficientes constantes}
Para EDOS com a cara
\[ a y''+ b y' + c y = 0 \]
onde \( a, b \)  e \( c \) são constantes, soluções podem sembre ser dadas
pelo chute \( y_1(x) = e^{r x} \) e a solução geral \( y(x) = C_1 e^{r_1 t} + C_2 e^{r_2 t}  \).

Substituindo na EDO genérica dada:
\[ a r^2 e^{rx} + b r e^{rx} + c e^{rx} = 0  \]
\[ e^{rx} \left( a r^2 + b r + c \right) = 0 \]

\[ \Leftrightarrow ar^2 + br + c = 0 \]


\begin{theorem}[Equação caractéristica de uma EDO]
  Dada uma EDO linear homogênea de segunda ordem com coeficientes constantes:
  $a y'' + b y' + cy = 0$, a equação \( a r^2 + br + c = 0 \) é a equação caractéristica
  da EDO dada.
\end{theorem}

Sejam \( r_1  \) e \( r_2 \) solução da equação caractéristica, temos três possíveis casos:

\paragraph{$\blacksquare$ Caso} \( r_1, r_2 \in \mathbb{R} \)  e \( r_1 \neq r_2 \): \( y(x) =  \alpha e^{r_1 x} + \beta e^{r_2 x}\) 
\paragraph{$\blacksquare$ Caso} \( r_1 = r_2 \), usar método do Wronskiano ou redução de ordem.
\paragraph{$\blacksquare$ Caso} \( r_1 \neq r_2 \)  e \( r_1 = \overline{r_2} \in \mathbb{C} \): \( y(x) = \alpha e^{r_1x} + \beta e^{r_2 x} \) 
  Nesse caso é necessário fazer uma mudança de coordenadas e passar a trabalhar com coordenadas
  polares e tratar usando formúla de Euler. Lembrando que as EDOs que estamos trabalhando estão em
  \( \mathbb{R} \), logo temos que sempre tomar cuidado em checar se a solução final está nesse domínio.

\subsection{Solução para coeficientes não constantes}
Equações com a cara 
\[ a(t) y''(t) + b(t) y'(t) + c(t) y(t) = f(t) \]

Solução geral: \( y(t) = y_{H} (t) + y_{P} (t) \) 

onde \( y_H \) é uma solução geral da EDO homogênea associada (obtida substituindo
$f(t)$ por zero e resolvendo a EDO homogênea resultante) e \( y_P \) é uma solução particular
qualquer da EDO inomogênea dada.

Supondo que já resolvemos a EDO homogênea por algum método:

\[ y_H(t) = \alpha  y_1 (t) + \beta y_2(t) \]
onde já conhecemos as funções \( y_1 \) e \( y_2 \) .


\subsection{Método da variação de parâmetros para encontrar \( y_P(t) \)}

Reescrevendo a solução particular da inomogênea como :

\[ y_P(t) = \alpha (t) y_1(t) + \beta (t) y_2(t) \]

E colocando a EDO original na forma canônica:
\[ y_P' = \alpha y_1 ' + \alpha ' y_1 + \beta  y_2 ' + \beta ' y_2 \]

Impomos a relação:
\begin{equation}
  \alpha ' y_1 + \beta ' y_2 = 0  
  \label{eq:inomogeneas_1}
\end{equation}

segue que \( y_P ' = \alpha (t) y_1 '(t) + \beta (t) y_2 ' (t) \) 


\[ y_P '' =  \alpha ' y_1 ' + \alpha y_1 ''  + \beta ' y_2 ' + \beta y_2 '' \]


Substituindo na EDO dada:
   
\[ y_P '' + p y_P ' + q y_P  = \alpha y_1 '' + \alpha ' y_1 ' + \beta y_2 '' + \beta ' y_2 ' + p
  \left( \alpha y_1 ' + \beta  y_2 '  \right) + q \left( \alpha y_1 + \beta  y_2 \right) = r(t) \]
\[ \alpha (t) \underbrace{\left[ y_1 '' (t) + p(t) y_1' (t)  + q(t) y_1(t) \right]}_{ =  0}  + \beta (t) \underbrace{\left[
      y_2''(t) + p(t) y_2 ' (t) + q(t) y_2 (t)  \right]}_{ = 0 } + \alpha ' y_1' + \beta ' y_2' = r(t) \]

\begin{equation}
  \alpha ' y_1 ' + \beta ' y_2 ' = r(t) 
  \label{eq:inomogeneas_2}
\end{equation}


Com as duas equações (\ref{eq:inomogeneas_1}) e (\ref{eq:inomogeneas_2}) temos o sistema linear:

\[ \begin{cases}
     y_1 \alpha ' + y_2 \beta ' = 0 \\
     y_1 ' \alpha ' + y_2 ' \beta ' = r(t)   
   \end{cases}\]
 
 \[ \begin{bmatrix}
      y_1 & y_2 \\ y_1 ' & y_2 '
    \end{bmatrix} =  \begin{bmatrix} \alpha ' \\ \beta ' \end{bmatrix}
    = \begin{bmatrix}  0 \\ r(t)    \end{bmatrix} \]
  
  \[  \alpha  ' = \frac{\det \begin{bmatrix} 0 && y_2 \\ r && y_2' \end{bmatrix}}
      {\det \begin{bmatrix} y_1 & y_2 \\ y_1 ' & y_2 ' \end{bmatrix}} = -
      \frac{r(t) y_2 (t) }{y_1 y_2' - y_2 y_1 } = - \frac{r(t) y_2(t) }{W[y_1, y_2](t)} \]
    \[ \beta  ' = \frac{\det \begin{bmatrix} y_1 & 0 \\ y_1 ' & r(t) \end{bmatrix} }
    {\det \begin{bmatrix} y_1 & y_2 \\ y_1 ' & y_2 ' \end{bmatrix}} = -
    \frac{r(t) y_1 (t) }{y_1 y_2' - y_2 y_1 } = - \frac{r(t) y_1(t) }{W[y_1, y_2](t)} \]

Integrando as soluções temos \( \alpha , \beta \) , após isso substituimos em \( y_P(t) = \alpha (t)
y_1(t) + \beta (t) y_2(t)\)  e a
solução geral da EDO será :

\[ y(t) = \left( \alpha _0 + \alpha (t) \right) y_1(t) + \left( \beta _0 + \beta (t) \right)y_2(t) \]

 \paragraph{$\blacksquare$ Exemplo}: \( x^2 y '' - 3x y' + 4y = x^2 \ln |x| , y(x) = ?, x> 0 \) 

 Considerando a EDO homogênea associada temos
 \[ x^2 y'' - 3xy' + 4y = 0  \]
 
 Suponto \( y(x) = x^{\lambda } \Rightarrow y' = \lambda x^{\lambda -1} \Rightarrow y'' = \lambda
 (\lambda - 1) x^{\lambda - 2}  \) 
 
 \[ x^{\lambda } \left[ \lambda (\lambda - 1) - 3 \lambda  + 4  \right] = 0 \Rightarrow \lambda
   ^2 - 4 \lambda  + 4 = 0 \Rightarrow \left( \lambda - 2 \right)^2 = 0 \Rightarrow \lambda  = 2 \]
 Logo, \( \lambda = 2 \) e uma das soluções da EDO homogenea associada é \( y_1 (x) = x^2 \) 


 Usando método do Wronskiano para encontrar \( y_2(t) \)

 \[ y_1 y_2 ' - y_1 ' y_2 = W_0 e^{- \int p(x) dx}, p(x) = - \frac{3}{x} \]
 \[ x^2 y_2 ' - 2x y_2 = e^{3 \ln x} = e^{x^3}  \]
 \[ y_2 ' - \frac{2}{x} y_2 = x \Rightarrow \mu = \frac{1}{x^2} \]
 \[ \left( \mu y_2 \right)' = \frac{1}{x} \Rightarrow \frac{y_2}{x^2} = \ln x + \mathbb{C}
   \Rightarrow y_2(x) = x^2 \ln x + x^2 \mathbb{C}\]
 Fazendo \( \mathbb{C} = 0 \), temos que \( y_2(x) = x^2 \ln x \) 
 uma possível solução é \( y_2(x) = x^2 \ln x \)   e a outra é \( y_1(x) = x^2 \) 
 Para encontrar \( y_P(x) \) impomos:
 \[ y_P(x) = \alpha (x) y_1(x) + \beta (x) y_2 (x) = \alpha (x) x^2 + \beta (x) x^2 \ln x\]
 \[ y_P'(x) =  \alpha ' x^2 + 2 \alpha  x + \beta ' x^2 \ln x + \beta (2x\lnx + x) \]
 impondo também \( \alpha ' x^2 + \beta ' x^2 \ln x = 0  \) 
 
 \[ y_P '   =  2 \alpha x + \beta (2x \ln x + x) \]
 \[ y_P'' = 2 \alpha ' x + 2 \alpha + \beta ' (2x \ln x + x ) + \beta (2 \ln x + 3) \]
 Substituindo na EDO:
 \[ y_P '' - \frac{3}{x} y_P ' + \frac{4}{x^2} y_P = \ln x \]
 \begin{multline*}
   \Leftrightarrow \left[ 2 \alpha ' x + 2 \alpha  + \beta ' (2x \ln x + x) + \beta (2 \ln x + 3)
   \right] - \frac{3}{x} \left[ 2 \alpha x + \beta \left( 2x \ln x + x \right) \right]+
   \\ \frac{4}{x^2} \left( \alpha x^2 + \beta x^2 \ln x  \right) = \ln x \end{multline*}
 \[ 2 \alpha ' x + \left( 2x \ln x + x \right)\beta ' = \ln x \]
 Temos então o sistema
 $$\begin{cases}
     x^2 \alpha ' + 2 x \ln x \beta ' = 0 \\
     2 \alpha ' x + \left( 2x \ln x + x \right) \beta ' = \ln x  
   \end{cases}$$
   
   \[ \Leftrightarrow \begin{cases}
          \alpha ' + \ln x \beta ' = 0 \\ 2 \alpha ' + \left( 2 \ln x + 1\right) \beta ' = \frac{\ln x}{x}
        \end{cases} \]

      \[ \alpha ' = - \beta ' \ln x   \]
      \[ - \beta ' \ln x + 2\ln x \beta' + \beta ' = \frac{\ln x}{x} \]
      \[ \beta ' = \frac{\ln x }{x} \Rightarrow \alpha ' = - \frac{(\ln x)^2}{x} \]

      \[ \beta = \int \frac{\ln x}{x} dx \text{ e } \alpha = - \int \frac{(\ln x)^2}{x} dx \]
      \[ \beta = \frac{(\ln x)^2}{2} \text{ e } \alpha = - \frac{(\ln x )^3}{3} \]

      \[ y_P(x) = - \frac{(\ln x )^3 x^2}{3} + \frac{(\ln x)^2 x^2 \ln x}{2}\]
      \[ y_P(x) = \frac{x^2 (\ln x)^3}{6} \]

      A solução geral é:

      \[ y(x) =  \alpha _0 x^2 + \beta _0 x^2 \ln x + \frac{x^2 (\ln x)^3 }{6}\]
      onde \( \alpha_0 \)   e \( \beta_0 \) são as condições de contorno.

      

      \[ a(t) y''(t) + b(t) y'(t) + c(t) y(t) = f(t) \]
      
      Solução geral: \( y(t) = y_{H} (t) + y_{P} (t) \) 

      onde \( y_H \) é uma solução geral da EDO homogênea associada (obtida substituindo
      f(t) por zero e resolvendo a EDO homogênea resultante) e \( y_P \) é uma solução particular
      qualquer da EDO inomogênea dada.
      
      Supondo que já resolvemos a EDO homogênea por algum método:
      
      \[ y_H(t) = \alpha  y_1 (t) + \beta y_2(t) \]
      onde já conhecemos as funções \( y_1 \) e \( y_2 \) .
      
      
      Método da variação de parâmetros para encontrar \( y_P(t) \) 

      Reescrevendo a solução particular da inomogênea como :

      \[ y_P(t) = \alpha (t) y_1(t) + \beta (t) y_2(t) \]
      
   E colocando a EDO original na forma canônica:
   \[ y_P' = \alpha y_1 ' + \alpha ' y_1 + \beta  y_2 ' + \beta ' y_2 \]

   Impomos a relação:
   \begin{equation}
   \alpha ' y_1 + \beta ' y_2 = 0  
   \label{eq:inomogeneas_1}
   \end{equation}
   segue que \( y_P ' = \alpha (t) y_1 '(t) + \beta (t) y_2 ' (t) \) 


   \[ y_P '' =  \alpha ' y_1 ' + \alpha y_1 ''  + \beta ' y_2 ' + \beta y_2 '' \]


   Substituindo na EDO dada:

   \[ y_P '' + p y_P ' + q y_P  = \alpha y_1 '' + \alpha ' y_1 ' + \beta y_2 '' + \beta ' y_2 ' + p
   \left( \alpha y_1 ' + \beta  y_2 '  \right) + q \left( \alpha y_1 + \beta  y_2 \right) = r(t) \]


   \[ \alpha (t) \underbrace{\left[ y_1 '' (t) + p(t) y_1' (t)  + q(t) y_1(t) \right]}_{ =  0}  + \beta (t) \underbrace{\left[
   y_2''(t) + p(t) y_2 ' (t) + q(t) y_2 (t)  \right]}_{ = 0 } + \alpha ' y_1' + \beta ' y_2' = r(t) \]

   \begin{equation}
   \alpha ' y_1 ' + \beta ' y_2 ' = r(t) 
   \label{inomogeneas_2}
   \end{equation}


   Com as duas equações (\ref{inomogeneas_1}) e (\ref{eq:inomogeneas_2}) temos o sistema linear:

   \[ \begin{cases}
   y_1 \alpha ' + y_2 \beta ' = 0 \\
   y_1 ' \alpha ' + y_2 ' \beta ' = r(t)   
    \end{cases}\]

   \[ \begin{bmatrix}
   y_1 & y_2 \\ y_1 ' & y_2 '
   \end{bmatrix} =  \begin{bmatrix} \alpha ' \\ \beta ' \end{bmatrix}
= \begin{bmatrix}  0 \\ r(t)    \end{bmatrix} \]

\[  \alpha  ' = \frac{\det \begin{bmatrix} 0 && y_2 \\ r && y_2' \end{bmatrix} }{\det \begin{bmatrix}
y_1 & y_2 \\ y_1 ' & y_2 ' \end{bmatrix}} = -
\frac{r(t) y_2 (t) }{y_1 y_2' - y_2 y_1 } = - \frac{r(t) y_2(t) }{W[y_1, y_2](t)} \]

\[  \beta  ' = \frac{\det \begin{bmatrix} y_1 & 0 \\ y_1 ' & r(t) \end{bmatrix} }{\det \begin{bmatrix}
y_1 & y_2 \\ y_1 ' & y_2 ' \end{bmatrix}} = -
\frac{r(t) y_1 (t) }{y_1 y_2' - y_2 y_1 } = - \frac{r(t) y_1(t) }{W[y_1, y_2](t)} \]


Integrando as soluções temos \( \alpha , \beta \) , após isso substituimos em \( y_P(t) = \alpha (t)
y_1(t) + \beta (t) y_2(t)\)  e a
solução geral da EDO será :


\begin{equation}
  y(t) = \left( \alpha _0 + \alpha (t) \right) y_1(t) + \left( \beta _0 + \beta (t) \right)y_2(t) 
  \label{eq:solucao_geral_inomogenea}
\end{equation}

 \paragraph{$\blacksquare$ Exemplo} : \( x^2 y '' - 3x y' + 4y = x^2 \ln |x| , y(x) = ?, x> 0 \) 

 Considerando a EDO homogênea associada temos
 \[ x^2 y'' - 3xy' + 4y = 0  \]

 Supondo \( y(x) = x^{\lambda } \Rightarrow y' = \lambda x^{\lambda -1} \Rightarrow y'' = \lambda
 (\lambda - 1) x^{\lambda - 2}  \) 
 
 \[ x^{\lambda } \left[ \lambda (\lambda - 1) - 3 \lambda  + 4  \right] = 0 \Rightarrow \lambda
   ^2 - 4 \lambda  + 4 = 0 \Rightarrow \left( \lambda - 2 \right)^2 = 0 \Rightarrow \lambda  = 2 \]
 Logo, \( \lambda = 2 \) e uma das soluções da EDO homogenea associada é \( y_1 (x) = x^2 \) 
 
 
 Usando método do Wronskiano para encontrar \( y_2(t) \)
 
 \[ y_1 y_2 ' - y_1 ' y_2 = W_0 e^{- \int p(x) dx}, p(x) = - \frac{3}{x} \]
 \[ x^2 y_2 ' - 2x y_2 = e^{3 \ln x} = e^{x^3}  \]
 \[ y_2 ' - \frac{2}{x} y_2 = x \Rightarrow \mu = \frac{1}{x^2} \]
 \[ \left( \mu y_2 \right)' = \frac{1}{x} \Rightarrow \frac{y_2}{x^2} = \ln x + \mathbb{C}
   \Rightarrow y_2(x) = x^2 \ln x + x^2 \mathbb{C}\]
 Fazendo \( \mathbb{C} = 0 \), temos que \( y_2(x) = x^2 \ln x \) 
 
 uma possível solução é \( y_2(x) = x^2 \ln x \)   e a outra é \( y_1(x) = x^2 \) 
 
 
 Para encontrar \( y_P(x) \) impomos:
 
 \[ y_P(x) = \alpha (x) y_1(x) + \beta (x) y_2 (x) = \alpha (x) x^2 + \beta (x) x^2 \ln x\]
 \[ y_P'(x) =  \alpha ' x^2 + 2 \alpha  x + \beta ' x^2 \ln x + \beta (2x\lnx + x) \]
 
 
 impondo também \( \alpha ' x^2 + \beta ' x^2 \ln x = 0  \) 
 
 \[ y_P '   =  2 \alpha x + \beta (2x \ln x + x) \]
 \[ y_P'' = 2 \alpha ' x + 2 \alpha + \beta ' (2x \ln x + x ) + \beta (2 \ln x + 3) \]
 
 Substituindo na EDO:
 \[ y_P '' - \frac{3}{x} y_P ' + \frac{4}{x^2} y_P = \ln x \]
 \begin{multline*}
   \Leftrightarrow \left[ 2 \alpha ' x + 2 \alpha  + \beta ' (2x \ln x + x) + \beta (2 \ln x + 3)
   \right] - \frac{3}{x} \left[ 2 \alpha x + \beta \left( 2x \ln x + x \right) \right]+
   \\ \frac{4}{x^2} \left( \alpha x^2 + \beta x^2 \ln x  \right) = \ln x \end{multline*}
 
 \[ 2 \alpha ' x + \left( 2x \ln x + x \right)\beta ' = \ln x \]
 
 
 
 Temos então o sistema
 
 $$\begin{cases}
     x^2 \alpha ' + 2 x \ln x \beta ' = 0 \\
     2 \alpha ' x + \left( 2x \ln x + x \right) \beta ' = \ln x  
   \end{cases}$$
   
   \[ \Leftrightarrow \begin{cases}
          \alpha ' + \ln x \beta ' = 0 \\ 2 \alpha ' + \left( 2 \ln x + 1\right) \beta ' = \frac{\ln x}{x}
        \end{cases} \]
      
      \[ \alpha ' = - \beta ' \ln x   \]
      \[ - \beta ' \ln x + 2\ln x \beta' + \beta ' = \frac{\ln x}{x} \]
      \[ \beta ' = \frac{\ln x }{x} \Rightarrow \alpha ' = - \frac{(\ln x)^2}{x} \]
      
      
      
      \[ \beta = \int \frac{\ln x}{x} dx \text{ e } \alpha = - \int \frac{(\ln x)^2}{x} dx \]
      
      \[ \beta = \frac{(\ln x)^2}{2} \text{ e } \alpha = - \frac{(\ln x )^3}{3} \]
      
      
      \[ y_P(x) = - \frac{(\ln x )^3 x^2}{3} + \frac{(\ln x)^2 x^2 \ln x}{2}\]
      \[ y_P(x) = \frac{x^2 (\ln x)^3}{6} \]
      
   
      
      A solução geral é:

      \[ y(x) =  \alpha _0 x^2 + \beta _0 x^2 \ln x + \frac{x^2 (\ln x)^3 }{6}\]
      onde \( \alpha_0 \)   e \( \beta_0 \) são as condições de contorno.



%%% Local Variables:
%%% mode: latex
%%% coding: utf-8
%%% TeX-master: "main.tex"
%%% End:



\chapter{Método de Séries}
\section{Mini revisão de séries de Taylor}
\subsection{Multiplicação de duas séries}
  Dadas as funões analiticas em \( x_0 \) , $f(x)$ e $g(x)$, com
  \[ f(x) = \sum_{n = 0}^{\infty} a_n (x - x_0)^n \text{ e } g(x) = \sum_{n = 0}^{\infty}b_n (x - x_0)^n \] 
  vale a seguinte relação
  \[ f(x) g(x) := \bigg[\sum_{n = 0}^{\infty} a_n (x - x_0)^n  \bigg] \bigg[\sum_{n = 0}^{\infty} b_n (x - x_0)^n  \bigg] \]

  \[ f(x) g(x) := \sum_{n = 0}^{\infty} c_n (x-x_0)^n \]
  onde \( c_n = a_0 b_n + a_1 b_{n-1} + \dots + a_n b_0 \) .

\subsection{Série de taylor em torno de um ponto \( x_0 \):}
\[ \boxed{ f(x) =  \sum_{n = 0}^{\infty} \frac{f^{(n)}(x_0)}{n!} (x-x_0)^n } \]
\section{Solução em torno de pontos regulares}
\subsection{Introdução e analiticidade de séries de potências}
 Dada a série de potências $ \sum a_n (x-x_0)^n$, se a série
 converge para valores de \( x \) num intervalo \( x_0 - \delta, x_0 + \delta \) , \( \delta > 0 \) ,
 então a série é analitica em \( x = x_0 \) .
 \paragraph{$\blacksquare$ Exemplo}: Dada a EDO \( a(x) y''(x) b(x)y'(x) + c(x)y(x) = 0 \)
 \[ y''(x) + p(x) y'(x) + q(x) y(x) = 0  \]
 

 \begin{enumerate}
  \item[] Escolher o ponto \( x_0 \) que vamos analisar.
  \item[] Estudar a analiticidade da série.
  \end{enumerate}

\paragraph{$\blacksquare$ Exemplo}:  \( (1 - x) y'' + y = 0  \)

  \[ y'' + \frac{1}{1 - x}y = 0, p(x) = \frac{1}{1 - x} = \sum_{n = 0}^{\infty} x^n \]

  Pelo critério da razão: \( \lim_{n \infty} |r_n| = |x| \), convergente no intervalo \( (-1, 1) \)
  e convergente no ponto \( 0 \) . Escrevendo a EDO em formato de séries:

  \[ \sum_{n = 0}^{\infty} (n + 2)(n+ 1) a_{n+2} x^n + \bigg[ \sum_{n = 0}^{\infty} x^n \bigg]
  \bigg[ \sum_{n = 0}^{\infty} a_n x^n \bigg] = 0\]

\subsection{Solução em séries em torno de pontos regulares}

 \paragraph{$\blacksquare$ Exemplo}: \( y''(x) + \sin(x) y(x) = 0 \) 

    Temos \( p(x) = 0 \) e \( q(x) = \sin(x) = \sum_{n = 0}^{\infty} (-1)^n \frac{x^{(2n +
    1)}}{(2n + 1)!} \), ambas funções são analiticas em torno de \( x_0
    \). Escolhendo
    \( x_0 = 0 \) temos:

    \[ \sum_{n = 0}^{\infty} (n + 2)(n+1) a_{n + 2} x^n + \left[ \sum_{n = 0}^{\infty} (-1)^n
    \frac{x^{(2n + 1)}}{(2n + 1)!} \right] \left[ \sum_{n = 0}^{\infty} a_n x^n \right] = 0\]

    Nesse tipo de resolução podemos definir uma ordem n e tentar encontrar o termo geral \( a_n \) e
    então retornar à série original de \( y(x) \) nas condições iniciais dadas.
    

    Abrindo os somatórios e agrupando os termos a partir das menores potências de \( x \) até \( x^3 \)  ;

    \[ (2 \cdot 1 a_2 + 3 \cdot 2 a_3 x + 4 \cdot 3 a_4 x^2 + 5 \cdot  4 x^3 + \cdots) +
    \left[ x - \frac{x^3}{3!} + \frac{x^5}{5!} - \cdots  \right] \left[ a_0 + a_1 x + a_2 x^2 + a_3
    x^3 + \cdots  \right]\]
    \[ 2 a_2 + (6 a_3 + a_0) x + (12 a_4 + a_1) x^2 + (20 a_5 + a_2 - \frac{a_0}{6})x^3 +
    \mathcal{O}(x^4) = 0 \]

    Relação de independência linear, ou seja, os coeficientes *devem* ser nulos.

    \[ \begin{cases}
    a_2 = 0 \\
    6a_3 + a_0 = 0 \\
    12 a_4 + a_1 = 0 \\
    \vdots 
    \end{cases}  \Rightarrow  \begin{align} a_3 = - \frac{a_0}{3} \\ a_4 = - \frac{a_1}{12} \\ a_5
    =  - \frac{a_0}{120}\end{align}\]
  
    Escrevendo \( y(x) \) em termo dos coeficientes encontrados: 

    \[ y(x) =  a_0 + a_1 x  - \frac{a_0}{6}x^3 - \frac{a_1}{12}x^4 + \frac{a_0}{120} x^5 + \cdots \]
    \[ y(x) = a_0 \left[ 1 - \frac{x^3}{3!} + \frac{x^5}{5!} + \cdots  \right] + a_1 \left[ x -
    \frac{x^4}{12} + \cdots  \right]\]

    Onde \( a_0 \) e \( a_1 \) são as condições de contorno.

    \paragraph{$\blacksquare$Exemplo}: \( y'' + k^2 y = 0, k \in \mathbb{R} \text{ e } y (x) = ?  \)
    Exemplo de uma EDO que podemos resolver sem o método de séries. Caso do movimento oscilatório.


    De forma direta, temos

    \[ \sum_{n = 0}^{\infty} \left[ (n+2) (n+1) a_{n + 2} x^n \right]  + k^2 \sum_{n = 0}^{\infty} a_n x^n = 0\]
    \[ \sum_{n = 0}^{\infty} \left[ (n+2)(n + 1) a_{n + 2} + k^2 a_n \right] x^n = 0  \]
    \[ a_{n + 2} = - k^2 \frac{a_n}{(n + 2)(n + 1)}, n = 0, 1, 2, ... \]

     $$ \begin{flalign*}
    a_2 & = - k^2 \frac{a_0}{2} &\\
    a_3 & = -k^2 \frac{a_1}{3 \cdot  2} &\\
    a_4 & = -k^2 \frac{a_2}{4 \cdot 3 } = k^4 \frac{a_0}{4 \cdot  3 \cdot  2 \cdot  1 }&\\
    a_5 & = -k^2 \frac{a_3}{5 \cdot 4 } = k^4 \frac{a_1}{5 \cdot  4 \cdot  3 \cdot  2 }
    \end{flalign*} $$

    Separando em termos pares e ímpares:

    \item{$\blacksquare$} n par: \( n = 2j, j = 0, 1, 2, ... \)
      \[ a_{2j + 2} = - k^2 \frac{a_{2j}}{(2j + 2)(2j + 1)} = - \frac{k^2}{(2j+2)(2j+1)} \left[ -
      \frac{k^2 a_{2j - 2}}{(2j)(2j -1)} \right]  \]
      \[ a_{2j + 2} = \frac{(-k^2)(-k^2)\cdot  \cdots (-k^2) a_0}{(2j+2)(2j + 1)(2j)(2j - 1) \cdots
      2)} = \frac{(-k^2)^{j + 1} a_0}{(2j + 2)!}\]
      \[ a_{2j + 2} = \frac{(-k^2)^{j + 1} a_0}{[2(j + 1)]!} \Rightarrow \boxed{ a_{2j} = (-1)^j\frac{
      k^{2j} a_0}{(2j)!}} \]
      
    \item{$\blacksquare$} n ímpar: \( n = 2j + 1 \) 

      \[ a_{2j + 3} = - \frac{k^2 a_{2j + 1}}{(2j + 3)(2j + 2)} = \frac{(-k^2)(-k^2) a_{2j-1}}{(2j+3)(2j+2)(2j+1)(2j)\cdots } \]
      \[ a_{2j + 3} =  \frac{(-k^2)(-k^2) \cdots (-k^2) a_1}{(2j+3)(2j+2)(2j+1)\cdots  3 \cdot  2} \]
      \[ a_{2j+3} = \frac{(-1)^j (k^2)^{j + 1} a_1}{(2j + 3)!} \Rightarrow \boxed{ a_{2j + 1} =
      (-1)^{j} \frac{(k)^{2j} a_1}{(2j + 1)!}} \]

    Voltando à série original, temos
    \[ y(x) = \sum_{n = 0}^{\infty} a_n x^n = \sum_{\text{n par}} a_n x^n + \sum_{\text{n impar}}
    a_n x^n \]
    
    \[ y(x) = a_0 \sum_{n = 0}^{\infty} (-1)^n \frac{k^{2n} x^{2n}}{(2n)!} + \underbrace{a_1 \sum_{n =
    0}^{\infty}(-1)^n \frac{k^{2n} x^{2n + 1}}{(2n + 1)!}}_{= \frac{a_1}{k} \left[ kx - \frac{k^3
    x^3}{3!} + \frac{k^5 x^5}{5!} - \frac{k^7 x^7}{7!} + \cdots  \right]} \]

    \[ y(x) = a_0 \cos(kx)  + \frac{a_1}{k} \sin(kx) \]

    \[ \boxed{ y(x) = A \cos(kx) + B \sin (kx) } \]

\section{Método de séries para pontos singulares regulares}



\[ y'' + p(x) y' + q(x) = 0 \]

Se \( p(x) \) e \( q(x) \) forem analíticas em \( x = x_0 \), então \( x = x_0 \) é um \textbf{ponto
ordinário} e vale a solução do tipo
\[ y(x) = \sum_{n = 0}^{\infty} a_n (x - x_0)^n \]
caso contrário, \( x = x_0 \) é um ponto singular, portanto precisamos de um outro tratamento
para poder escrever a solução no formato de série de potências.

Se \( p(x) \) \textbf{ou} \( q(x) \) não forem analíticas em \( x = x_0 \) , mas \( (x - x_0) p(x) \) e \(
(x-x_0)^2 q(x) \) forem, então
\( x = x_0 \) é dito ser um *ponto singular regular*. Caso contrário \( x = x_0 \) é um ponto
singular irregular \emph{(esse caso não é tratado nessa disciplina, então não tentei buscar entender ainda...)}.

  \paragraph{$\blacksquare$ Exemplo}: \( x^2 y'' + 2x y' + 3y = 0 \Leftrightarrow y'' + \frac{2}{x}y ' + \frac{3}{x^2}y = 0 \)

  \[ xp(x) = 2 \text{ e } x^2 q(x) = 3 \]

  No caso \( x = x_0 \) ser ponto singular regular, procuramos uma solução em série de
  potências com a seguinte forma:


  \begin{equation}
    y(x) = (x - x_0)^{\lambda} \sum_{n = 0}^{\infty} a_n (x - x_0)^n, a_0 \neq 0
    \label{eq:solucao_singular_regular}
  \end{equation}


  Impomos que ao menos uma das soluções devem ser dessa forma e caso o \( \lambda = 0  \), ou seja,
  conseguimos apenas uma solução a partir dele, podemos chegar à outras soluções por outros métodos
  como redução de ordem ou Wronksiano...

  Expandindo \( (x - x_0) p(x) \) e \( (x - x_0)^2 q(x) \) em séries de potências:


  $$ \begin{cases}
       (x - x_0) p(x) = \sum_{n = 0}^{\infty} p_n (x - x_0)^n \\
       (x - x_0)^2 q(x) = \sum_{n = 0}^{\infty} q_n (x - x_0)^n
     \end{cases} $$
     
     
     \[ y(x) = \sum_{n = 0}^{\infty} a_n (x - x_0)^{n + \lambda}\]
     \[ y'(x) = \sum_{n = 0}^{\infty} (n + \lambda)a_n (x-x_0)^{n + \lambda - 1} \]
     \[ y''(x) = \sum_{n = 0}^{\infty} (n+\lambda)(n + \lambda - 1) a_n (x - x_0)^{n + \lambda - 2} \]
     
     Saindo da forma canônica multiplicando a EDO por (x - x_0)^2:

     \[ (x - x_0)^2 y'' + (x-x_0) \left[ (x-x_0) p(x) \right]y' + \left[ (x-x_0)^2 q(x) \right]y = 0  \]
     
     \begin{multline*}
       (x-x_0)^2 \sum_{n = 0}^{\infty} (n+\lambda)(n + \lambda - 1) a_n x^{n + \lambda - 2} + (x-x_0)
       \sum_{n = 0}^{\infty} \left[ (x - x_0) p(x) \right] \left( (n + \lambda)a_n(x-x_0)^{n + \lambda - 1}
       \right) \\ + \sum_{n = 0}^{\infty} \left[ (x-x_0)^2 q(x) \right] a_n (x - x_0)^{n + \lambda} =  0
     \end{multline*}
     
     Fazendo a expansão em série de Taylor para \( (x - x_0) p(x)  \) e \( (x - x_0)^2 q(x) \) de forma a
     considerar apenas o termo mais baixo \( p_0, q_0 \), temos que
     \[ (x - x_0)p(x) = p_0 + (x-x_0) p_1 + (x - x_0)^2 p_2 \cdots \Rightarrow p_0 = \lim_{x\to x_0}
       (x-x_0) p(x)  \]
     \[ q_0 = \lim_{x\to x_0} (x-x_0)^2 q(x) \]
     
     
     Rearranjando as séries acima e simplificando as potências temos a seguinte série:
     
     $$ \sum_{n = 0}^{\infty} \left[ (n + \lambda) (n + \lambda - 1)a_n + [(x - x_0) p(x)] (n +
       \lambda)a_n + [(x-x_0)^2 q(x)] a_n \right] \left(x - x_0 \right)^{n + \lambda}=0 $$

     $$
     \sum_{n = 0}^{\infty} \left[ (n + \lambda)(n + \lambda - 1) + \left[ (x - x_0)p(x) \right] (n +
       \lambda) + \left[ (x-x_0)^2 q(x) \right] \right] a_n (x - x_0)^{n + \lambda} = 0
     $$


     \[  \sum_{n = 0}^{\infty} \left[ (n + \lambda )(n + \lambda - 1) + (n + \lambda ) p_0 + q_0 \right]
       a_n (x - x_0)^{n + \lambda} \]
     
     Abrimos as séries e escrevendo as combinações em termos de ordem mais baixa. Isso é útil porque
     assim
     podemos fatorar e chagar numa relação para \( \lambda  \):
     
     \[ \lambda (\lambda -1) a_0 \left( x - x_0 \right)^{\lambda } + p_0 \lambda a_0 (x -
       x_0)^{\lambda} + q_0 a_0 (x - x_0)^{\lambda } + \mathcal{O} (x - x_0)^{\lambda + 1} =  \]
     \[ a_0 \left[ \lambda (\lambda - 1) + \lambda p_0 + q_0  \right] \left( x - x_0 \right)^{\lambda } +
       \mathcal{O}(x - x_0)}^{\lambda + 1} = 0\]
   
   Podemos escrever a relação: 
   
   \begin{equation}
     \lambda \left( \lambda - 1 \right) + \lambda p_0 + q_0 = 0
     \label{eq:equacao_indicial}
   \end{equation}
   
   
   \paragraph{$\blacksquare$ Exemplo} : \( y'' + \frac{2}{x}y' + \frac{3}{x^2}y = 0, x_0 = 0 \) 
   
     \[ p(x) = \frac{2}{x} \Rightarrow xp(x) = 2 = p_0 + p_1 x + p_2 x^2 + ... \]
     
     \[ q(x) = \frac{3}{x^2} \Rightarrow x^2 q(x) = 3 = q_0 + q_1 x + q_2 x^2 + ... \]
     
     Logo, \( x_0 = 0 \) é um ponto singular regular e \( y(x) = x^\lambda \sum_{n = 0}^{\infty} a_n x^{n} \) .
     
     Encontrando o coeficiente \( \lambda  \) usando a equação indicial (\ref{eq:equacao_indicial}) :
     \[ \lambda (\lambda - 1) + p_0 \lambda  + q_0 = 0 \Rightarrow \lambda ^2 + \lambda  + 3 = 0 \]
     \[ \lambda  = \frac{-1 \pm \sqrt{-11}}{2}  \Rightarrow \lambda _1 = - \frac{1}{2} \pm
       \frac{i\sqrt{11}}{2}, \lambda _2 = - \frac{1}{2} - \frac{i\sqrt{11}}{2}\]
     
     \[ y(x) = y_1(x) + y_2(x) = x^{\lambda _1} \sum_{n = 0} a_n x^n + x^{\lambda _2} \sum_{n =
         0}^{\infty} a_n x^n \]
     \subsection{Casos para \( \lambda  \)}
     Esse indice é definido a partir da relação (\ref{eq:equacao_indicial}) e para cada solução da
     equação quadrática tem significados diferentes. Eles foram divididos em alguns casos.

     Nessa etapa do curso muitos desses resultados não foram rigorosamente justificados e
     em alguns deles não fui atrás de tentar entender como chegar nessas relações por
     conta própria. :p 😛
     
     \subsubsection{Caso \( \lambda  \pm \notin \mathbb{R}   ( \lambda _{-} = \lambda _{+}^{*}) \)}
     A solução de \( \lambda  \) é um complexo. Nesse caso o complexo conjugado também é válido, mas
     como queremos apenas soluções reais \( y(x) \) , precisamos tratar o caso complexo
     a fim de conseguir apenas a perte real da solução.
     Para isso utilizamos as seguintes soluções:
     
     \[ y_1 (x) = | x - x_0 |^{\lambda _+} \sum_{n = 0}^{\infty} a_n (x - x_0)^n, a_0 \neq 0 \]
     \[ y_2 (x) = | x - x_0|^{\lambda _{-}} \sum_{n = 0}^{\infty} b_n (x - x_0)^{n}, b_0 \neq 0 \]
     
     
     \subsubsection{Caso \( \lambda  \pm \in \mathbb{R} ( \lambda _+ \geq \lambda _-) \) }
     
     
     \begin{enumerate}
     \item[$\blacksquare$] Subcaso: \( y_1 (x) = |x - x_0 |^{\lambda _+} \sum_{n = 0}^{\infty} a_n (x - x_0)^n \)
       
       Nesse caso a segunda solução pode ser construida a partir da primeira e
       fica em termos de logaritmo e outra série. Só utilizamos o resultado.
       Se em algum momento eu estudar as demonstrações posso coloca-las aqui.

       \[ y_1 (x) = y_1(x) \ln |x - x_0| + |x - x_0|^{\lambda _+} \sum_{n = 0}^{\infty} b_n (x - x_0)^n \]
       
       Além disso, para definir os coeficientes \( b_n \) é necessário substitutir a solução na EDO
       original e encontrar o termo geral. Ao menos foi isso que entendi.
     
     
     \item[$\blacksquare$] Subcaso:\( \lambda _+ - \lambda _{-} \in \mathbb{N}^{*} \)       
       
       \[ y_1(x) = |x - x_0|^{\lambda } \sum_{n = 0}^{\infty} a_n (x - x_0)^n \]
       \[ y_2 (x) =  \alpha y_1(x) \ln |x - x_0| + |x - x_0|^{\lambda} \sum_{n = 0}^{\infty} b_n (x - x_0)^n \]
       
 
     \item[$\blacksquare$] Subcaso:  \( 0 < \lambda _{+} - \lambda _{-} \notin \mathbb{N} \)
       
       
       \[ y_1(x) =  | x - x_0|^{\lambda _+} \sum_{n=  0}^{\infty} a_n (x - x_0)^n\]
       \[  y_2 ( x) =  |x - x_0|^{\lambda_ - } \sum_{n = 0}^{\infty} b_n (x - x_0)}^n \]
     
   \end{enumerate}

   
   \paragraph{$\blacksquare$Exemplo}:   Equação de Bessel
   \[ x^2 y'' + x y' + (x^2 +  \nu^2) y = 0; \nu \in \mathbb{R}_{+} \]
   
   Para \( x_0 = 0  \) 
 
   \[  y'' + \frac{1}{x} y' + \left( 1 - \frac{\nu^2}{x^2} \right) y = 0 \]
   \[ p(x) = \frac{1}{x} \]e \( q(x) = 1 - \frac{\nu^2}{x 2} \) 
   
   Não são analiticas em \( x_0 \) , encontrando os pontos singulares regulares:

   \[ x p(x) = 1 \text{  e } x^2 q(x) = x^2 - \nu^2 \]
   
   essas funções são analiticas em \( x_0 \) .
   Resolvendo a equação indicial:
   
   \[ \lambda (\lambda - 1) p_0 \lambda + q_0 = 0 \]
   \[ p_0 =  \lim_{x\to x_0} xp(x) = 1 \text{ e }  q_0 = \lim_{x \to x_0} x^2q(x) = - \nu^2\]
   
   \[ \lambda (\lambda - 1) + \lambda - \nu^2 = 0  \]
   \[ \lambda ^2 - \nu^2 = 0 \Rightarrow \lambda  = \pm \nu  \]
   (Se o \( \mu \) não tivesse sido definido como estritamente positivo a solução seria \( \lambda  =
   \pm | \nu|\)).
   

   \paragraph{$\blacksquare$ Exercicio proposto}:  resolver para \( \nu = 0, \nu = \frac{1}{2}, \nu = 1 \) 
   
 
%%% Local Variables:
%%% mode: latex
%%% coding: utf-8
%%% TeX-master: "main.tex"
%%% End:


  
\chapter{Transformada de Laplace}
% falar sobre o que é uma transformada integral.
% Exemplo e definiçã 
A transformada de Laplace é uma transformada integral que associa funções $\mathbb{R} \mapsto \mathbb{C}$.
É usada para modelar sistemas em que há mudanças rápidas de comportamentos (a.k.a ``impulsos''). Em
geral comportamentos que saem de um regime homogêneo em um instante $t$.

\begin{theorem}
  Definição: No espaço das funções $\mathcal{F}$ temos funções do tipo $f:\mathbb{R}_{+} \mapsto \mathbb{C}$,
  definimos a transformada de laplace como
  \begin{equation}
  F(s) \equiv \mathcal{L} \{ f \} (s) = \int_{-\infty}^{\infty} e^{-st} f(t) dt
  \label{eq:transformada_laplace}
  \end{equation}
\end{theorem}

\begin{figure}[h!]
  \centering
  \includegraphics[width=0.6\textwidth]{./images/transformada_laplace.pdf}
\end{figure}


% Define the functions
\def\fx{exp(\x) / (\x + 1)}
\def\gx{(\x + 2) / (\x - 1)}
\def\hx{(\x - 4) / (\x + 3) + \x / 2 + 2}
\begin{figure}
  \centering
\begin{tikzpicture}
  % Axis
  \draw[->] (0,0) -- (7,0) node[right] {$x$};
  \draw[->] (0,-1) -- (0,7) node[above] {$y$};

  % Plot f(x) for 0 <= x <= 2
  \draw[domain=0:2, samples=100, smooth, red!80!black, thick] plot (\x, { \fx });

  % Plot g(x) for 2 <= x <= 4
  \draw[domain=2:4, samples=100, smooth, blue, thick] plot (\x, { \gx });

  % Plot h(x) for 4 <= x <= 6
  \draw[domain=4:6, samples=100, smooth, green!70!black, thick] plot (\x, { \hx });

  % Dashed vertical lines
  \draw[dashed] (2, 0) -- (2, 6);
  \draw[dashed] (4, 0) -- (4, 6);
  \draw[dashed] (6, 0) -- (6, 6);
\end{tikzpicture}
\end{figure}





\section{Propriedades necessárias de $f$}

\paragraph{(1) Seccionalidade contínua} [contínua por pedaços]


% IMAGEM


Note que nos extremos dos intervalos a função deve convergir para algum valor.


\paragraph{(2) Limite superior} Cresce mais devagar que ao menos \textbf{uma} exponencial.

Essa condição é similar à um limite superior/inferior da função.

$$ |f(t) | < M e^{at} , t > t_0, a \in \mathbb{R}_+  \Longrightarrow  \exists F(s), s > a  $$

% IMAGEM

Garantidamente a função estará contida na faixa destacada.




\paragraph{$\blacksquare$ Exemplo}: $f(t) = 1, t \geq 0$
$$ \mathcal{L} \{ 1 \} (s)  = \int_{0}^{\infty} e^{-st} dt = \left( - \frac{1}{s} \right) e^{-st} \bigg|_{0}^{\infty} =
\frac{1}{s}, s > 0$$
\paragraph{$\blacksquare$ Exemplo}: $f(t) = t$
$$ \mathcal\{ t \} (s) = \int_{0}^{-\infty} e^{-st} t dt = - \underbrace{\frac{e^{-st}}{s} t
  \bigg | _{0}^{\infty}}_{= 0} + \int_{0}^{\infty}
\frac{e^{-st}}{s} dt  $$
$$  \mathcal\{ t \} (s) = \int_{0}^{\infty} \frac{e^{-st}}{s} dt = \frac{1}{s^2} , s > 0 $$
\paragraph{$\blacksquare$ Exemplo}: $f(t) = t^n, n \in \mathbb{N}$
$$ \mathcal{L} \{ t^n \} (s) = \int_{0}^{\infty} e^{-st} t^n dt = \left( - \frac{1}{s} e^{-st} t^n \right)
\bigg|_{0}^{\infty} - \frac{n}{s} \int_{0}^{\infty} e^{-st} t^{n - 1} dt  $$

$$ \Leftrightarrow  \mathcal{L} \{ t^n \} (s) = \frac{n}{s} \mathcal{L} \{t^{n-1}} (s) = \frac{n}{s} \frac{(n - 1)}{s} \mathcal{L} \{t^{n-2}\} (s) $$
$$ \Leftrightarrow \mathcal{L} \{ t^n \} (s) = \frac{n(n-1) \cdot \cdots \cdot 1}{s \cdot s \cdots s} \mathcal{L} \{ 1 \} (s) = \frac{n!}{s^n} \mathcal{L} \{1\} (s) $$
$$ \Longrightarrow \mathcal{L} \{ t^n \} (s) = \frac{n!}{s^{n + 1}}  $$

\section{Propriedades da Transformada de Laplace}
\subsection {Propriedade da exponencial}
\paragraph{$\blacksquare$ Exemplo} $f(t) = e^{at}, t \geq 0$
$$F(s) = \int_{0}^{\infty} e^{-st} f(t) dt = \int_{0}^{\infty} e^{-st} e^{at} dt = \int_{0}^{\infty} e^{-(s-a)t} dt = -
\frac{e^{-(s-a)}}{s - a} \bigg |_{0}^{\infty}$$

$$ \Longrightarrow F(s) = \frac{1}{s-a}, s > 0 $$

\paragraph{$\blacksquare$ Exemplo} $\mathcal{L} \{1\} (s) = \frac{1}{s}, s> 0$
$$ \mathcal{L} \{ e^{at} \} (s) = \frac{1}{s - a}, s-a > 0 \Rightarrow s > a $$
\paragraph{$\blacksquare$ Exemplo} $g(t) = e^{-3t} t^4$
Sabendo que $\mathcal{L} \{ t^n \} (s) = \frac{n!}{s^{n+1}}$, temos 
$$ \mathcal{L} \{e^{-3t} t^4\} (s) = \mathcal{L} \{ t^4 \} (s + 3) = \frac{4!}{(s+3)^5}, s > 3  $$

\subsection{propriedade de translação}
A função $g(t)$ é dada por $g(t) = u_{c}(t) f(t - c)$ onde $u_{c}$ é uma função
degrau unitário.

\subsubsection{Função degrau unitário ou função de Heaviside ($\Theta(t - a)$)}
\begin{equation}
  u_{a} (t) = \begin{cases} 1, \text{se } t > a \\ 0, \text{se } t < a \end{cases}
\end{equation}


Sabendo a $\mathcal{L}\{f\} (s)$, podemos saber $\mathcal{L} \{g\}(s)$:

$$ G(s) = \int_{0}^{\infty} e^{-st} g(t) dt = \int_{0}^{\infty} e^{-st} u_c(t) f(t -c) dt $$
$$ G(s) = \int_{C}^{\infty} e^{-st} f(t) \overset{t' = t-c}{=} \int_{0}^{\infty} e^{-s(t'+c)} f(t') dt' $$
$$ G(s) = e^{-st} F(s) , C > 0 $$

\item[$\blacksquare$] Se $g(t) = f(ct), c > 0 \Rightarrow G(s) = \frac{1}{C} F(\frac{s}{c})$.
\item[$\blacksquare$] Se $g(t) = \frac{df}{dt}$, então $G(s) = \int_{0}^{\infty} e^{-st} g(t) dt = \int_{0}^{\infty} e^{-st}
  \left( \frac{df}{dt} \right) dt $

  $$ \Leftrightarrow G(s) = e^{-st} f(t) \bigg|_{0}^{\infty} - \int_{0}^{\infty} (-s)e^{-st} f(t) dt  $$
  $$ \Leftrightarrow G(s) = -f(0) + s \int_{0}^{\infty} e^{-st} f(t) dt   $$
  $$ \Leftrightarrow G(s) = -f_0 + s F(s)  $$
  $$ \Longrightarrow \mathcal{L} \bigg\{ \frac{df}{dt} \bigg\} (s) = -f_0 + s \mathcal{L} \{f\} (s)  $$

  Portanto, $\mathcal{L}\{f''\} (s) = -f'(0) + s \mathcal{L}\{f'\} (s)$

  $$ \mathcal{L}\{f''\} (s) = -f'(0) + s \left(-f_0 + s \mathcal{L}\{f\} (s) \right) $$
  $$ \mathcal{L} \{f''\} (s) = -f_0' - sf_0 + s^2 \mathcal{L} \{f\} (s) $$
  De forma geral, vale
  \begin{equation}
    \mathcal{L} \{ f^{(n)} \} (s) = -f(t_0)^{n-1} - s f(t_0)^{n-2} + \dots + s^{n} \mathcal{L} \{ f\} (s)
    \label{eq:transformada_derivada}
  \end{equation}

  Se $g(t) = t^n f(t), n \in \mathbb{N}$ sabemos qual é a transformada $\mathcal{L} \{f\} (s)$, então valoe
  $$ G(s) = (-1)^{n} F^{(n)} (s) \equiv (-1) \frac{d^{n}}{ds^{n}} F(s)  $$

  \paragraph{$\blacksquare$ Exemplo}

  $$ \mathcal{L} \{ 1 \} (s) = \frac{1}{s} \Rightarrow \mathcal{L} \{ t^n \} (s) = (-1)^{n} \frac{d^n}{ds^n} \left( \frac{1}{s}
  \right) = \frac{n!}{s^{n+1}}$$



  

\subsection{Delta de Dirac (Função impulso)}


[IMAGEM]

\( \delta _{t_0} (t) = 0 \) para \( t \neq t_0 \) e \( \int_{t_0 - \epsilon_1}^{t_0 + \epsilon_2} \delta_{t_0} (t) dt = 1 \) .
Note que o intervalo da função não precisa necessariamente ser simétrico, ou seja, podemos
ter \( |\epsilon_1| \neq  | \epsilon_2 | \) .


[IMAGEM]

Com \( a \mapsto 0 \) a função fica mais centrada em torno de \( t_0 \) .

\paragraph{Propriedade mais útil do Delta de Dirac}

\[ \int_{t_0- \epsilon_1}^{t_0 + \epsilon_2} \delta_{t_0} (t) f(t) dt = f(t_0), \epsilon_1, \epsilon_2 > 0 \]



\[ \mathcal{L}\{ \delta_{t_0}(t) \} (s) = \int_{0}^{\infty} e^{-st} \delta_{t_0}(t) dt = e^{-st_0}, t_0 > 0 \]

\[ \mathcal{L}\{ u_a(t) \} (s) = \frac{e^{-as}}{s} , a> 0\]


Para \( a > 0 \): \( \mathcal{L} \{ \delta_a (t)  \} (s) = s \mathcal{L}\{ u_a(t) \} (s) = -u_0 (0) + s \mathcal{L}\{ u_a(t) \} (s) \) 

ou seja, vale a relação \( \delta_a (t)\) ``$=$'' \( \frac{d}{dt}u_a(t) \) 


\section{Aplicação em EDOs: usando Transformada de Laplace em PVIs}

Nosso objetivo é usar a transformada de Laplace para resolver problemas do tipo

\[ a y'' + by' + cy = f(t), y(0) = y_0, y'(0) = v_0  \]

Podemos aplicar a propriedade (\ref{eq:transformada_derivada}) à EDOs e como
é uma transformada integral, sabemos que a transformação é linear, logo,
podemos var a operação em toda equação de forma que


\[ a\mathcal{L} \{y''\} (s) + b \mathcal{L} \{y'\} (s) + c \mathcal{L} \{y\} (s) = \mathcal{L} \{f\} (s) \]

\[ a \left( -v_0 - y_0 s + s^2 Y(S) \right)  + b \left( -y_0 sY(s) \right) + c Y(s) = F(s) \]

\[ \Leftrightarrow (a s^2 + bs + c) Y(s) = F(s) + (a v_0 + b y_0) + a y_0 s \]
\[ \Leftrightarrow Y(s) = \frac{F(s)}{(as^2 + bs + c)} + \frac{(a v_0 + b y_0 + a y_0 s)}{(as^2 + bs + c)} \]


\paragraph{$\blacksquare$ Exemplo}: \( y'' + w_0^2 y = \sin( \omega t), t > 0, y(0) = 0, y'(0) = 0 \) 
\[ - y'(0) - y(0) s + s^2 Y(s) + \omega_0 ^2 Y(s) = \mathcal{L} \{\sin(\omega t)\}  \]

\item[$\blacksquare$] Fazendo a transformada de Laplace da função \( \sin(\omega t)  \) :

  Trabalhar com número complexo pode ser mais simples já que no fim das contas vamos
  estar manipulando uma exponencial. É importante notar que como estamos querendo a
  transforamda para uma função em \( \mathbb{R} \), o resultado não pode ser em \( \mathbb{C} \),
  então é importante checar em qual domínio está a função final.

  \[ \mathcal{L} \{\sin (\omega t) = \mathcal{L} \bigg\{ \frac{e^{i\omega t} - e^{-i \omega t} }{2i}\bigg\} \}(s) = \frac{1}{2i} \mathcal{L}
    \{e^{i\omega t}\} - \frac{1}{2i} \mathcal{L} \{e^{-i\omega t}\}   \]

  \[ \mathcal{L} \{\sin (\omega t) \}   = \frac{1}{2i}\frac{1}{(s - i \omega)} - \frac{1}{2i} \frac{1}{(s + i \omega)} =
    \frac{s + i \omega - s - i \omega}{2i (s-i\omega)(s + i \omega)} \]

  \[ \mathcal{L} \{\sin(\omega t)\} (s) = \frac{\omega}{s^2 + \omega^2}  \]
  
  

  Voltando à resolução da EDO:


  \[ \Leftrightarrow s^2 Y(s) + \omega_0^2 Y(s) = \frac{\omega}{s^2 + \omega^2}   \]


  \[ Y(s) = \frac{\omega}{(s^2 + \omega_0^2)(s^2 + \omega^2}\]
  \item[$\blacksquare$] Tratando a fração parcial:

    \[  \frac{\omega}{(s^2 + \omega_0^2) (s^2 + \omega^2) }  = \frac{A}{(s^2 + \omega^2} + \frac{B}{s^2 + \omega_0^2} \]

    \[ \Leftrightarrow \frac{A(s^2 + \omega_0^2) + B(s^2 + \omega^2)}{(s^2 + \omega^2)(s^2 + \omega_0^2)} \Rightarrow \begin{cases}
                                                                            A + B = 0 \Rightarrow B = - A \\
                                                                            A\omega_0^2 - A \omega^2 = \omega \Rightarrow A =
                                                                            \frac{\omega}{\omega_0^2 - \omega^2}
                                                                          \end{cases}  \]

    \[ \Leftrightarrow  Y(s) = \frac{\omega}{\omega_0^2 - \omega^2} {\left[ \frac{1}{s^2 + \omega^2} - \frac{1}{s^2 +
            \omega_0^2} \right]}\]                                                                      

    A partir daqui adaptamos os termos dentro dos colchetes para chegar em uma relação que fique
    claro alguma transformada de Laplace conhecida, nesse caso uma \( \mathcal{L} \{\sin\}  \) .

    \[ \frac{1}{s^2 + \omega^2} = \frac{1}{\omega} \mathcal{L} \{\sin(\omega t)\} (s) = \frac{1}{\omega} \frac{\omega}{s^2 + \omega^2} \]
    \[ \frac{1}{s^2 + \omega_0^2} = \frac{1}{\omega_0} \mathcal{L} \{\sin(\omega_0 t)\} (s) = \frac{1}{\omega_0}\frac{\omega_0}{s^2 + \omega^2} \]


    \[ Y(s) = \frac{\omega}{\omega_0^2 - \omega^2} \left[ \frac{1}{\omega} \mathcal{L} \{\sin(\omega t) \} - \frac{1}{\omega_0} \mathcal{L} \{\sin(\omega_0
        t)\}  \right] \]


    \[ Y(s) = \mathcal{L} \bigg\{ \frac{\sin(\omega t)}{\omega_0^2 - \omega^2}  - \frac{(\omega / \omega_0) \sin(\omega_0 t)}{\omega_0^2 - \omega^2}\bigg\}(s)  \]

    \[ \Leftrightarrow Y(s) = \mathcal{L} \{y(t)\} (s) \Leftrightarrowq y(s) = \frac{\sin(\omega t) - (\omega/\omega_0) \sin(\omega_0 t)}{\omega_0^2
        - \omega^2} \]
    com \( \omega_0^2 \neq \omega^2 \) .

\section{Teorema da convolução e transformada de Laplace Inversa}


\begin{theorem}
  \textbf{Definição}: produto de convolução para transformada de Laplace.
  Sejam duas funções \( f(t) \) e \( g(t) \), o produto de convolução delas é definido por
  \begin{equation}
    \left( f * g \right)(t) := \int_{0}^{t} f(u) g(t-u) du \equiv \int_{0}^{t} g(u) f(t-u)du = \left( g * f
    \right) (t)
  \end{equation}
\end{theorem}

vale a propriedade de assortatividade,  \( \left( f * 0 \right)(t) = 0 \)  e \(  \left( f * 1
\right)(t) \neq 1 \).


Se\( \mathcal{L} \{f\}  = F \)  e \( \mathcal{L} \{g\}  = G \), então

\begin{equation}
\mathcal{L} \{f * g\} (s) = F(s) G(s)  
\end{equation}



\item[$\blacksquare$] Pequeno esboço do caminho que deve seguir para demonstrar esse resultado:

  \[ F(s) G(s) =  \left(  \int_{0}^{\infty} e^{-st} f(t) dt  \right) \left( \int_{0}^{\infty} e^{-st'} g(t')dt'
    \right) = \int_{0}{\infty} \int_{0}^{\infty} e^{-st} f(t) e^{-st'} g(t') dt' dt \]
  
  \[ F(s) G(s) =  \int_{0}^{\infty} \int_{0}^{\infty} e^{-s(t+t')} f(t) g(t') dt' dt \]

  fazendo \( t = u - t' \) e ajustando os limites de integração e o produto de convolução dentro da
  integral externa.


  \paragraph{$\blacksquare$ Exemplo}: Seja \( H(s) = \frac{1}{s^2(s^2 + 4)} = F(s)G(s) \) 

  \[ \mathcal{L} \{t\} (s) = \frac{1}{s^2} = F(s) \]

  \[ \mathcal{L} \{\frac{\sin(2t)}{2}\} (s) = \frac{1}{s^2 + 4} = G(s) \]

  \[ \Rightarrow h(s) = \mathcal{L} \{f * g\} (s) = \int_{0}^{t} g(u)f(t-u) du = \int_{0}^{t} \frac{\sin(2t)}{2} (t - u)du \]

  \[ h(s) = \frac{t}{2} \int_{0}^{t} \sin(2u)du - \frac{1}{2} \int_{0}^{t} \sin(2u) u du \]

  \[ \Rightarrow h(s) = -\frac{t}{4} \left[ \cos(2u) \right]_{0}^{t} - \frac{1}{2} \bigg\{ \left[ -
      \frac{\cos(2u)}{2}u \right]_{0}^{t} + \int_{0}^{t} \frac{\cos(2u)}{2} du  \bigg\} \]

  \[ h(s) = - \frac{t}{4} \left[ \cos(2t) - 1 \right] + \frac{\cos(2u)}{4}- \frac{\sin(2t)}{8} \]

  \[ h(s) = - \frac{\sin(2t)}{8} + \frac{t}{4} \]



%%% Local Variables:
%%% mode: latex
%%% coding: utf-8
%%% TeX-master: "main.tex"
%%% End:


% \chapter{Extra - Alguns exercicios/exemplos interessantes}
% \paragraph{Resolvendo circuito RLC}

\chapter*{Apêndice}
\section{Método de resolução de integral de fração parcial}
\paragraph{Exemplo} \( f(x) = \frac{1}{ax^2 + bx + c} \Rightarrow \int f(x) dx \) 

Encontrando as raízes do polinômio (e/ou simplificando)
\[ ax^2 + bx + c = 0 = a(x-r_1)(x_r_2) \]
onde \( r_1 \) e \( r_2 \) são raízes da equação quadrática.
Escrevendo em outra forma:
\[ \frac{1}{a(x-r_1)(x-r_2)} = \frac{A}{x-r_1} + \frac{B}{x-r_2} = \frac{A(x-r_2) +
    B(x-r_1)}{(x-r_1)(x-r_2)} = \frac{(A+B)x - (Ar_2 + Br_1)}{(x-r_1)(x-r_2)} \]
como não há termo de \( x \)  no numerador \( A + B = 0 \Rightarrow B = -A \) , então
\[ \frac{1}{a(x-r_1)(x- r_2)} = - \frac{A(r_2 - r_1)}{(x-r_1)(x - r_2)} \]
\[ -A \underbrace{(r_2 - r_1)}_{= \sqrt{\Delta}} = \frac{1}{a} \Rightarrow A = - \frac{1}{\sqrt{\Delta}}\]

\[ \int \frac{dx}{ax^2 + bx + c} = - \frac{1}{\sqrt{\Delta}} \bigg\{ \int \frac{dx}{x - r_1} - \int \frac{dx}{x - r_2} \bigg\} \]
\[ = - \frac{1}{\sqrt{\Delta}} \left( \ln |x - r_1| + \ln | x - r_2 | \right) = \frac{1}{\sqrt{\Delta}} \ln
  \bigg|\frac{x - r_1}{x - r_2} \bigg| \]
% fazer um exemplo de casos em que a EDO não é tão simples
% e falar sobre a relação (x^3 + x^2 + ...) = (Ax + B)/...


\nocite{*}
\printbibliography
\end{document}
%%% Local Variables:
%%% mode: latex
%%% TeX-master: "main.tex"
%%% End:
